%% ---------------------------------------------------------------------------
%% What binds a surrogate — a pre-registered falsification ladder
%% Target: Computer Methods in Applied Mechanics and Engineering (Elsevier)
%%
%% EVERY NUMBER IN THIS FILE IS A MACRO RESOLVED FROM A RESULTS FILE.
%% Build with:  python paper/numbers.py && python build_paper.py
%% A bare numeral in the prose is a build failure, not a style issue: defects
%% B21 and B43 both began with a number that reached a document with no script
%% behind it. See paper/numbers.py for the file and path behind each macro.
%% ---------------------------------------------------------------------------
\documentclass[preprint,11pt]{elsarticle}

\usepackage{amsmath,amssymb}
\usepackage{booktabs}
\usepackage{enumitem}
\usepackage{graphicx}
\usepackage{siunitx}
\usepackage[hidelinks]{hyperref}
\usepackage{xcolor}

% GENERATED BY paper/numbers.py — DO NOT EDIT.
% Every value below is read from a results file; see SPEC in numbers.py
% for the file and path behind each one.
\newcommand{\nLzeroTracer}{1.69e-03}
\newcommand{\nLzeroTracerFront}{0.0}
\newcommand{\nLzeroTracerPe}{93}
\newcommand{\nLzeroBCdiff}{3.5}
\newcommand{\nLzeroRetarded}{0.0}
\newcommand{\nLzeroRetardedR}{901}
\newcommand{\nLzeroThermal}{0.0}
\newcommand{\nLzeroLDF}{3.0e-04}
\newcommand{\nIsoNmat}{240}
\newcommand{\nIsoKHmin}{0.0322}
\newcommand{\nIsoKHmed}{0.359}
\newcommand{\nIsoKHmax}{2.52}
\newcommand{\nIsoKHdec}{1.9}
\newcommand{\nIsoHenryMed}{-2.0}
\newcommand{\nIsoHenryMax}{-0.9}
\newcommand{\nIsoHenryBelowMilli}{73}
\newcommand{\nIsoHenryBelowMicro}{34}
\newcommand{\nAlphaFolds}{0.02}
\newcommand{\nNclustFolds}{240}
\newcommand{\nFprFoldsNominal}{6.5}
\newcommand{\nFprFoldsCal}{2.8}
\newcommand{\nAlphaManifest}{0.02}
\newcommand{\nNclustManifest}{48}
\newcommand{\nFprLegacyNominal}{8.5}
\newcommand{\nNclustLegacy}{12}
\newcommand{\nLoneMLP}{0.0306}
\newcommand{\nLoneXGB}{0.0361}
\newcommand{\nLoneRF}{0.0388}
\newcommand{\nLoneRidge}{0.0516}
\newcommand{\nLonePodFloor}{2.80e-04}
\newcommand{\nLoneMLPvsXGBdiff}{-0.00548}
\newcommand{\nLoneMLPvsXGBlo}{-0.00675}
\newcommand{\nLoneMLPvsXGBhi}{-0.00426}
\newcommand{\nLCbeta}{0.222}
\newcommand{\nLCbetaLo}{0.205}
\newcommand{\nLCbetaHi}{0.239}
\newcommand{\nLCsmallest}{0.0654}
\newcommand{\nLClargest}{0.0348}
\newcommand{\nLClargestFixed}{0.0348}
\newcommand{\nLCsmallestFixed}{0.0657}
\newcommand{\nLCratio}{1.88}
\newcommand{\nLClastDiff}{0.00534}
\newcommand{\nLClastLo}{0.00439}
\newcommand{\nLClastHi}{0.00637}
\newcommand{\nLCmodesSmall}{2}
\newcommand{\nLCmodesLarge}{13}
\newcommand{\nLCcondBeta}{0.122}
\newcommand{\nLtwoBest}{0.0241}
\newcommand{\nLtwoBestFam}{mlp_depth}
\newcommand{\nLtwoBestCfg}{d8}
\newcommand{\nLtwoPodFloor}{2.80e-04}
\newcommand{\nLtwoDepthShape}{U-SHAPED, optimum at d8}
\newcommand{\nLtwoWidthShape}{U-SHAPED, optimum at w64}
\newcommand{\nLtwoXgbShape}{U-SHAPED, optimum at md6}
\newcommand{\nLtwoRfShape}{falls monotonically with capacity}
\newcommand{\nLtwoDepthMDE}{5.0}
\newcommand{\nLtwoXgbMDE}{2.0}
\newcommand{\nLthreeMlpSmall}{0.0564}
\newcommand{\nLthreeMlpLarge}{0.0572}
\newcommand{\nLthreeRbfLarge}{0.0702}
\newcommand{\nLthreeChebyLarge}{0.1028}
\newcommand{\nLthreeMlpLr}{3e-05}
\newcommand{\nLthreeChebyLr}{0.03}
\newcommand{\nLthreeArms}{63}
\newcommand{\nLthreeBestKanMlp}{0.0573}
\newcommand{\nLthreeBestKan}{0.0633}
\newcommand{\nLthreeBestKanDiff}{-0.0060}
\newcommand{\nLthreeBestKanLo}{-0.0120}
\newcommand{\nLthreeBestKanHi}{-0.0008}
\newcommand{\nLfiveFNO}{0.0216}
\newcommand{\nLfiveFNOmodes}{4}
\newcommand{\nLfiveFNOwidth}{28}
\newcommand{\nLfiveONet}{0.0265}
\newcommand{\nLfiveOKAN}{0.0357}
\newcommand{\nLfiveFloor}{5.02e-04}
\newcommand{\nLfiveONetFNOdiff}{0.00488}
\newcommand{\nLfiveONetFNOlo}{0.00035}
\newcommand{\nLfiveONetFNOhi}{0.00978}
\newcommand{\nLfiveFlatDiff}{-0.00102}
\newcommand{\nLfiveFlatLo}{-0.00585}
\newcommand{\nLfiveFlatHi}{0.00491}
\newcommand{\nLfiveFloorSmall}{1.44e-02}
\newcommand{\nLfiveFloorLarge}{5.02e-04}
\newcommand{\nLfiveBasisSmall}{0.01440}
\newcommand{\nLfiveBasisLarge}{0.00050}
\newcommand{\nLfiveCoefSmall}{0.0515}
\newcommand{\nLfiveCoefLarge}{0.0510}
\newcommand{\nLsixJoint}{0.0562}
\newcommand{\nLsixSeparate}{0.0541}
\newcommand{\nLsixDiff}{0.00210}
\newcommand{\nLsixLo}{0.00056}
\newcommand{\nLsixHi}{0.00363}
\newcommand{\nLsixSlope}{0.00140}
\newcommand{\nLsixSlopeLo}{-0.00329}
\newcommand{\nLsixSlopeHi}{0.00618}
\newcommand{\nLsixSlopeR}{0.040}
\newcommand{\nLsixBetweenSD}{0.0104}
\newcommand{\nLsixMDE}{4.0}
\newcommand{\nLsixPowerFour}{82}
\newcommand{\nLsixPowerTwo}{26}
\newcommand{\nLsixSizeNull}{1.8}
\newcommand{\nLsixBetweenOverBase}{18.5}
\newcommand{\nLsixDaMin}{7.7}
\newcommand{\nLsixDaMax}{173}
\newcommand{\nLsixDaMedian}{25.0}
\newcommand{\nLsixDaDecades}{1.35}
\newcommand{\nLsixKineticRtwo}{-0.256}
\newcommand{\nLsixStepRtwo}{0.98}
\newcommand{\nLsixQmaxRtwo}{0.93}
\newcommand{\nLsevenLearned}{0.0474}
\newcommand{\nLsevenKlinkenberg}{0.1769}
\newcommand{\nLsevenShock}{0.3146}
\newcommand{\nLsevenPattern}{0.3878}
\newcommand{\nLsevenDiff}{-0.1295}
\newcommand{\nLsevenLo}{-0.1384}
\newcommand{\nLsevenHi}{-0.1203}
\newcommand{\nLsevenRatio}{3.73}
\newcommand{\nWarpFixed}{0.05005}
\newcommand{\nWarpPredicted}{0.04770}
\newcommand{\nWarpOracle}{0.02203}
\newcommand{\nWarpPredDiff}{0.00235}
\newcommand{\nWarpPredLo}{-0.00812}
\newcommand{\nWarpPredHi}{0.01183}
\newcommand{\nWarpOracleDiff}{0.02802}
\newcommand{\nWarpOracleLo}{0.01990}
\newcommand{\nWarpOracleHi}{0.03668}
\newcommand{\nWarpHeadroom}{2.17}
\newcommand{\nNwidthOriginal}{31}
\newcommand{\nNwidthSingle}{102}
\newcommand{\nNwidthTwoWaveA}{17}
\newcommand{\nNwidthTwoWaveB}{13}
\newcommand{\nGapMin}{0.025}
\newcommand{\nGapMax}{0.992}
\newcommand{\nMonoRawNrmse}{0.04770}
\newcommand{\nMonoProjNrmse}{0.04812}
\newcommand{\nNwidthExpC}{-2.45}
\newcommand{\nNwidthRtwoC}{0.9910}
\newcommand{\nAnchorTfiftyExp}{287}
\newcommand{\nAnchorTfiftyModel}{303}
\newcommand{\nAnchorPlateauExp}{0.267}
\newcommand{\nAnchorPlateauModel}{0.336}
\newcommand{\nAnchorTfiveExp}{86}
\newcommand{\nAnchorTfiveModel}{1.7}
\newcommand{\nAnchorNrmse}{0.128}
\newcommand{\nGates}{23}
\newcommand{\nMassClosure}{0.0}
\newcommand{\nMassClosureMof}{0.1}
\newcommand{\nGridConv}{0.1}
\newcommand{\nGridConvMof}{0.3}
\newcommand{\nGridNz}{2000}
\newcommand{\nLoneXfloor}{109}
\newcommand{\nLtwoXfloor}{86}
\newcommand{\nLtwoVsLonePct}{21.3}
\newcommand{\nLtwoVsLoneLo}{18.7}
\newcommand{\nLtwoVsLoneHi}{23.7}
\newcommand{\nLtwoIdentity}{0.00e+00}
\newcommand{\nAnchorNrmseComsol}{0.069}
\newcommand{\nLsixLegacyBetweenSD}{0.0391}
\newcommand{\nLsixLegacyBase}{0.0566}
\newcommand{\nLfourbVerdictMaterial}{\textbf{[PENDING]}}
\newcommand{\nLfourbVerdictTime}{\textbf{[PENDING]}}
% derived ratios — the formula behind each is in DERIVED in numbers.py
% FloorDrop: POD floor at p=8 divided by the floor at p=128
\newcommand{\nFloorDrop}{28.7}
% FNOxFloor: the best FNO divided by the POD floor at p=128
\newcommand{\nFNOxFloor}{43}
% ONetxFloor: the best DeepONet divided by the POD floor at p=128
\newcommand{\nONetxFloor}{53}
% FNOgain: the best DeepONet divided by the best FNO
\newcommand{\nFNOgain}{1.23}
% NwidthHalving: modes needed to halve the error, 2^(1/|algebraic exponent|)
\newcommand{\nNwidthHalving}{1.76}
% LsixNoiseDrop: the legacy between-material sd divided by v2's
\newcommand{\nLsixNoiseDrop}{3.8}
% LsixLegacyBetweenOverBase: the legacy between-material sd as a share of its base error
\newcommand{\nLsixLegacyBetweenOverBase}{69}
% GapRange: the outlet separation of the two waves, max divided by min
\newcommand{\nGapRange}{40}
% LCmaterialsForHalving: materials needed to halve the error at the measured exponent
\newcommand{\nLCmaterialsForHalving}{23}


\journal{Computer Methods in Applied Mechanics and Engineering}

\begin{document}

\begin{frontmatter}

\title{What binds a surrogate: a pre-registered falsification ladder for
transfer to unseen adsorbent parameter sets in metal--organic framework water-adsorption
columns}

\author{Abdulhamid Batayhi}
\address{Department of Mechatronics Engineering, Faculty of Mechanical Engineering,
Y\i{}ld\i{}z Technical University, Istanbul, T\"urkiye\\
ORCID: \href{https://orcid.org/0009-0007-0374-6755}{0009-0007-0374-6755}}
\ead{abdulbatayhi@gmail.com}

\begin{abstract}
We ask what \emph{binds} the error of adsorption-column surrogates on adsorbents they
have never seen. Seven candidate explanations were named in advance, each with the
evidence that would eliminate it, and tested on a solver verified against
\nLzeroNchecks{} closed-form solutions and compared without fitting to a published
MOF-303 breakthrough curve. Each verdict is reported in pre-declared words, with the
design it rests on (\nNclustFolds\ materials in five folds, a split of
\nLfourbMatNmat, or a legacy set of \nNclustLegacy). For an adsorbent parameter set it
has never seen, a perceptron trained on the largest set predicts the breakthrough field to
\nLClargestPct\,\% normalised error. Training materials bind: error still falls as
\(N_{\mathrm{mat}}^{-\nLCbeta}\) at the largest size affordable. Capacity is not
eliminated. A perceptron beats Kolmogorov--Arnold networks at matched size in every
comparison, \nLthreeSigSim\ of \nLthreePairs\ significantly once tuning choices are
paid for (legacy set). The error is limited by how well the model learns the link
from isotherm and pellet properties to the field, not by its ability to represent
the field: improving the representation \(\nLfiveVtwoFloorDrop\times\) changes the
error by \nLfiveVtwoCoefChangePct\,\%. The learner matters (legacy set: a Fourier
neural operator is \(\nOperatorLever\times\) more accurate than the best pointwise
regressor, DeepONet \(\nONetLever\times\)), but none approaches what the
representation allows. Decomposing equilibrium and kinetics gives a small real
benefit with no detectable Damk\"ohler dependence; un-fitted closed forms are
\nLsevenRatio\ times worse than learning; aligning the two adsorption fronts would pay
\nWarpHeadroomFixed{}\(\times\) only if the fronts were known (legacy set). The
physics-residual rung is withdrawn and being re-run: its residual moved the gas at
the wrong velocity. \nLedgerA\ withdrawn claims and \nLedgerB\ caught defects are
reported in full.
\end{abstract}

\begin{keyword}
surrogate models \sep model order reduction \sep adsorption column dynamics \sep
transfer to unseen adsorbent parameter sets \sep Kolmogorov \(n\)-width \sep pre-registration
\end{keyword}

\end{frontmatter}

%% ===========================================================================
\section{Introduction}
\label{sec:intro}

The case for a learned surrogate of a fixed-bed adsorption column is a screening
case. A designer choosing among candidate metal--organic frameworks does not want
one breakthrough curve faster; they want thousands of them, for materials that have
not been simulated. The quantity that matters is therefore not fit but
\emph{transfer}: the error on a material the model has never seen.

That framing is unusual enough to be worth stating plainly, because it changes what
counts as evidence. A surrogate cannot beat its own reference on accuracy --- the
reference defines the target, and Section~\ref{sec:reference} measures rather than
assumes how good it is --- so a synthetic study can only measure
accuracy per unit inference cost, or accuracy from less information. Two defects in
our own frozen protocol came from forgetting this, and both are recorded in
Section~\ref{sec:corrections}.

It also changes what counts as a strong result. McGreivy and Hakim \cite{mcgreivy2024weak}, surveying
machine learning for fluid-related partial differential equations, determine that
\emph{79\,\% (60 of 76)} of articles claiming to outperform a standard numerical
method compare against a weak baseline, and find reporting biases widespread. That
finding is the design brief for this study rather than a citation of convenience.
Grossmann and co-workers \cite{grossmann2024pinn}, comparing physics-informed networks against the finite
element method on Poisson, Allen--Cahn and semilinear Schr\"odinger problems, report
that the networks did not outperform FEM \emph{in their study} --- a scoped result,
and we scope ours the same way.

\paragraph{Terms used throughout.} The paper uses a small vocabulary from
reduced-order modelling and from pre-registered testing, defined here once for
readers from reticular chemistry and adsorption engineering
(Figure~\ref{fig:schematic} shows the objects involved).
\begin{description}[leftmargin=1.2em,itemsep=0.1em,font=\mdseries\itshape]
\item[Surrogate.] A fast learned model that maps a parameter set describing an
adsorbent and an operating condition to the space--time concentration, loading and
temperature fields the column simulation would produce.
\item[Rung.] One pre-registered candidate explanation of the surrogate's error,
tested by one experiment whose possible outcomes and their wording were fixed before
it ran. The seven rungs form the \emph{ladder}.
\item[Binds.] An intervention binds the error if improving it improves the error on
held-out adsorbents; a rung is \emph{eliminated} when the pre-declared test shows it
does not.
\item[Basis and map.] A field is written as a weighted sum of fixed spatial--temporal
patterns, the \emph{basis} (proper orthogonal decomposition); the surrogate's real job
is the \emph{map} from the parameter set to the weights. The basis floor is the error
left when the weights are exact.
\item[\(n\)-width.] How fast the best possible basis error falls as patterns are added;
a theoretical lower bound for any method that sums a fixed basis.
\item[Twin.] The data-only model a physics-informed arm is compared with: identical
architecture, data and budget, without the physics term.
\item[Damk\"ohler number.] The product of the linear-driving-force rate constant and
the simulated duration (Section~\ref{sec:methods}); large values mean the solid stays
near equilibrium over the run, small values kinetic control. Because some runs were
extended, it also carries the length of the run.
\item[Henry branch, step.] The low-humidity, near-linear part of the water isotherm,
and the sharp cooperative step of pore filling above it; together they produce the two
adsorption fronts of Section~\ref{sec:physics}.
\item[Edge rule.] An optimum found at the edge of a tuning range does not count until
the range is extended past it, unless the metric has demonstrably stopped changing.
\end{description}

\begin{figure}[t]
\centering
\includegraphics[width=\textwidth]{../figures/Fig0_schematic.pdf}
\caption{What a surrogate maps. \textbf{A}~A packed bed of adsorbent pellets fed humid
air; the eleven numbers of Table~\ref{tab:params} describe the adsorbent, the bed and
the operating condition. \textbf{B}~One stored solver field of the larger dataset,
\(c/c_{\mathrm{in}}\) over position and time, chosen by a fixed rule (the clearest
complete two-front breakthrough). \textbf{C}~Its outlet curve: a fast front on the
Henry branch, a slow rise, and the sharp cooperative pore-filling front. A surrogate
must reproduce such fields for an adsorbent it has never seen.}
\label{fig:schematic}
\end{figure}

\paragraph{What this paper is not.} It is not an architecture paper. Rung L3
refutes the premise that a richer per-edge basis helps here, on our own data; rung
L5 refutes our own explanation of why the operators fail. The negative results are
the contribution.

\paragraph{Prior art, stated up front.} Transfer to held-out \emph{systems} is
established in operator learning: Poseidon \cite{herde2024poseidon} pretrains on fluid-dynamics equations and
reports that it ``generalizes very well to new physics that is not seen during
pretraining'' across fifteen downstream tasks, and Multiple Physics Pretraining \cite{mccabe2023multiple}
reports transfer to ``systems with previously unseen physical components''.
Held-out \emph{materials} is established for interatomic potentials: MOFSimBench \cite{krass2025mofsimbench}
evaluates twenty universal machine-learning interatomic potentials on
metal--organic frameworks. In adsorption-process engineering a surrogate that takes
the material as an input is also established: MAPLE \cite{pai2020generalized} takes
Langmuir isotherm parameters as inputs and predicts cyclic-steady-state performance
indicators for an arbitrary adsorbent, and it has been validated on a laboratory rig
with two zeolites whose measured isotherms were not in its training data
\cite{pai2022experimental}. Machine-learning surrogates of pressure-swing adsorption
have been used for multiobjective process optimisation \cite{subraveti2019machine},
and adsorbents, metal--organic frameworks among them, have been screened at the
process level with the help of learned models
\cite{leperi2019development,burns2020prediction}. On the numerical side, the arms of
our first rungs are a non-intrusive reduced-order model in the sense of Hesthaven and
Ubbiali \cite{hesthaven2018non}: a proper-orthogonal-decomposition basis with a
learned map from parameters to its coefficients. To the surrogate, a ``material'' is a
point in the space of isotherm and bed parameters of Table~\ref{tab:params}, with its
kinetics derived from them, and a sample adds the operating conditions; the
question this paper asks --- which term binds the error, the basis or the map --- is
the parametric reduced-order-model question, asked of an adsorption column and
answered with a pre-registered test. Operator learning has reached cyclic adsorption too, with
generalisation tested across initial conditions \cite{ceccanti2026deep}. What we do
not find is the combination tested here: transient fields rather than steady-state
indicators, inflected water isotherms rather than Langmuir ones, and a held-out
material test that is cluster-robust and calibrated and asks what binds the error
rather than whether it is small. We find no operator-learning study that evaluates
transfer to held-out materials --- a claim we
state with the papers that would otherwise falsify it cited alongside it, because an
earlier and broader version of it was falsified by exactly those papers
(Section~\ref{sec:corrections}, A23).

Heinlein and Taraz \cite{heinlein2026error} establish, for DeepONets on Korteweg--de~Vries and Burgers
equations, that the error is branch-dominated at large internal dimension, that the
learned trunk can be replaced by a classical basis, and that the branch exhibits
mode-wise spectral bias. We reproduce that signature independently and extend it in
four ways: the inputs are a finite-dimensional parameter vector describing a
\emph{material} rather than a discretised input function; evaluation is transfer to
held-out materials under a cluster-robust, calibrated test rather than
in-distribution error; we show the plateau is invariant to basis size while the
proper-orthogonal-decomposition floor beneath it falls \(\nFloorDrop\times\), and calibrate
the residual against an oracle to an equivalent mode count; and we locate it in the
map from material parameters to the field rather than in the basis. A similar plateau on a different input type and a
different partial differential equation suggests the wall is not peculiar to either
problem, though one study on each side is not a general law. A separate line of work formalises how a neural
operator relates to its own discretisation \cite{bartolucci2023reno};
our encoding is held fixed across every arm, so that axis is controlled rather
than studied here.

\paragraph{Contributions.}
\begin{enumerate}
\item A pre-registered falsification ladder over seven named candidate explanations
for the transfer error of an adsorption-column surrogate, with the eliminating
evidence declared before each run and every null reported with its minimum
detectable effect.
\item A reference solver verified against \nLzeroNchecks{} closed-form solutions
(Section~\ref{sec:reference}) and compared, with no parameter fitted to the curve,
against a published MOF-303 packed bed whose two-wave shape and half-breakthrough time
it reproduces and whose first-wave timing it does not (Section~\ref{sec:anchor}) --- a
single shallow bed, which is an anchor and not a validation.
\item A direct measurement of what binds the error --- the parameter-to-coefficient
map, not the representation --- in the \emph{optimal} basis, so the conclusion
cannot depend on the basis any particular method learned
(Section~\ref{sec:mechanism}).
\item A measurement of what an established change of coordinates buys on this
system, and of exactly what consumes it (Section~\ref{sec:warp}).
\item A correction record: \textbf{\nLedgerA} withdrawn claims and \textbf{\nLedgerB} caught
defects, reported as Section~\ref{sec:corrections} rather than hidden.
\end{enumerate}

%% ===========================================================================
\section{The reference and its verification}
\label{sec:reference}

Every surrogate number in this paper is \emph{surrogate-versus-solver}. The solver
is therefore not a detail, and it is verified rather than asserted.

The model is a one-dimensional non-isothermal packed bed with axial dispersion, a
linear-driving-force kinetic closure and a coupled energy balance. The inlet is a
third- or flux-type boundary condition. That choice is load-bearing and was
initially made wrongly: our first verification compared the solver against Ogata and
Banks \cite{ogata1961solution}, who solve the \emph{first-} or concentration-type inlet. Against the correct
closed form --- van Genuchten and Alves \cite{vangenuchten1982analytical}, who give the third-type solution --- the
error falls to \nLzeroTracer{} absolute, an improvement of more than an order of
magnitude, with the residual maximum at the outlet where a finite column legitimately
differs from a semi-infinite solution. The two boundary conditions differ by
\nLzeroBCdiff\,\% of the inlet concentration at the front half-height at
\(\mathrm{Pe} \approx \nLzeroTracerPe\); the difference does not stay confined to an
inlet boundary layer, which is why using the wrong reference is not a small matter.

The same choice is vindicated independently. Van Genuchten and Alves state that the
flux-type condition conserves mass inside a column whereas the concentration-type
condition causes mass-balance errors that grow with \(D/v\). That is precisely the
defect we found numerically and separately: a Dirichlet ghost cell in the dispersion
stencil drove an artificial dispersive influx and left global mass closure off by
about one per cent, invariant across snapshot counts, until the inlet dispersive flux was
set to zero, after which closure improved to \(\nMassClosure\,\%\)
(Section~\ref{sec:corrections}, B8).

\nLzeroNchecks{} closed-form solutions, each isolating one term:

\begin{table}[t]
\centering
\caption{Verification of the reference solver. Each check isolates one term of the
model against a closed form; \(R_f\) is the retardation factor of the front.
Figure~\ref{fig:ladder}C.}
\label{tab:l0}
\resizebox{\textwidth}{!}{%
\begin{tabular}{llr}
\toprule
check & term isolated & error \\
\midrule
inert tracer vs.\ van Genuchten third-type & advection, dispersion, the inlet & \nLzeroTracer{} abs \\
retarded front, \(R_f = \nLzeroRetardedR\) & isotherm coupling, equilibrium limit & \nLzeroRetarded\,\% \\
thermal wave, adsorption off & the energy equation & \nLzeroThermal\,\% \\
linear driving force vs.\ Anzelius--Schumann \cite{anzelius1926uber,schumann1929heat} & kinetics & \nLzeroLDF{} abs \\
\bottomrule
\end{tabular}}
\end{table}

Global mass balance closes to \(\nMassClosure\,\%\); grid convergence is \(\nGridConv\,\%\) at
\(N_z = \nGridNz\) for the generic configuration and \(\nGridConvMof\,\%\) for the
MOF-303-parameterised one.

\paragraph{What we read, and what we read through.} The Anzelius and Schumann papers
are closed and carry no abstract in any bibliographic service, so the \(J\)-function
was verified through Perry's handbook \cite{levan2019adsorption} rather than from either original.
That is a two-step chain, and we state both steps: Perry's carries the integral
definition of \(J\) and itself attributes it to a later chemical-engineering source we
could not verify, so our check rests on Perry's transcription. The same is
true of Danckwerts's 1953 paper \cite{danckwerts1953continuous}, for which we
additionally cite Pearson's note \cite{pearson1959note}.
Section~\ref{sec:reproducibility} lists every reference cited through a secondary
source rather than read at the primary. A study that audits its own numbers should
audit its own reading.

%% ===========================================================================
\section{The physics of the column}
\label{sec:physics}

\subsection{The isotherm must be inflected, with a finite Henry slope}

Water uptake in atmospheric-water-harvesting frameworks is small below a threshold
relative humidity --- in our parameter space the primary sites hold a fraction
\(f_H = \nDcFhLo\)--\(\nDcFhHi\) of capacity there (Table~\ref{tab:params}) --- then
rises almost vertically as pore filling proceeds cooperatively, then saturates. That step is the entire reason these materials work.
Single-site Langmuir cannot produce it: its second derivative is negative
everywhere, verified numerically as zero inflection points, and it predicts the
largest uptake gradient as \(c \to 0\) --- the opposite of the mechanism. A bare
Hill form produces the step but gives \(q \propto c^{n}\) at the origin, so
\(\mathrm{d}q/\mathrm{d}c \to 0\) and Henry's law is violated.

We therefore use a dual-term form \emph{in the spirit of} Do and Do \cite{dodo2000model}: a Langmuir
primary-site term carrying the Henry limit, plus a cooperative Sips term producing
the step,
\begin{equation}
q^{*}(c,T) = q_{\max}\left[
f_H\,\frac{b_H c}{1 + b_H c}
+ (1-f_H)\,\frac{(b_C c)^{n}}{1 + (b_C c)^{n}} \right],
\label{eq:isotherm}
\end{equation}
with \(b_H(T) = b_{H0}\exp(-\Delta H/RT)\) and \(b_C(T) = b_{C0}\exp(-\Delta H/RT)\).
Both affinities carry the same \(\Delta H\), so \(q_{st} = -\Delta H + RT\) holds at
every loading; it is reproduced to within \nIsoCCmaxErr\,\% at a quarter, a half and three quarters of capacity.

The phrase ``in the spirit of'' is deliberate. Do and Do superpose an \(n\)-layer
BET primary term with a Sips term; ours is a Langmuir primary term with a Sips term.
What we take from them is the two-term construction and the fact that \emph{only the
primary term carries the Henry slope} --- their Sips term, like ours, has slope
exactly zero at the origin \cite{buttersack2019modeling}. We read the Do--Do functional form itself through Buttersack's transcription, because the original paper could not be retrieved; Buttersack proposes a competing isotherm, so his transcription is not an endorsement of the form. Their cluster exponent is fixed; ours is a sampled
material parameter, which is the later model's treatment \cite{dodo2009new}. Calling equation~\eqref{eq:isotherm} ``the Do--Do form'' would
be wrong, and it is a correction we made to our own text
(Section~\ref{sec:corrections}, B56). We call the isotherm family S-shaped, or
stepped, and attach no IUPAC type to it here. The lower end of the sampled
cooperativity is only weakly stepped: \nChemWeakNPct\,\% of the materials of the
larger dataset have \(n < \nChemWeakN\), where the cooperative term is close to
Langmuir.

\subsection{Henry's law holds everywhere; the region where it dominates does not}

A property verified at one configuration is not a property of a parameter space ---
this is the error class of A18, A21 and A26 applied to the isotherm rather than to a
sample size. Measured across all \nIsoNmat\ sampled materials rather than at two
named configurations: \(n > 1\) strictly for every material, so the cooperative term
is \(o(c)\) at the origin and the Henry constant is finite and positive everywhere.
But \(K_H\) spans \nIsoKHmin--\nIsoKHmax\ \si{\mole\per\kilo\gram} per
\si{\mole\per\metre\cubed}, \nIsoKHdec\ decades, median \nIsoKHmed. And the
concentration at which the cooperative term reaches 1\,\% of loading --- the top of
the accessible Henry region --- reaches \(10^{\nIsoHenryMax}\) of the feed at best
and a median of \(10^{\nIsoHenryMed}\), but lies below \(10^{-3}\) of the feed for
\nIsoHenryBelowMilli\ materials and below \(10^{-6}\) for \nIsoHenryBelowMicro. For
those the isotherm is effectively Sips-like at every concentration the column
visits. In the units of a water isotherm: at fixed temperature the concentration is
proportional to relative humidity, so for a median material the near-linear uptake
regime ends at \(10^{\nIsoHenryMed}\) of the feed humidity, and for
\nIsoHenryBelowMicro\ materials below a millionth of it --- in practice, no linear
region at all. The design claim is thermodynamic validity by construction, which holds
everywhere; it is not an experimentally accessible linear regime, which does not.

\subsection{An inflected isotherm produces a two-wave breakthrough}

The consequence is that breakthrough is not a sigmoid. A fast, spreading wave on the
Henry branch arrives first and plateaus; a slow, self-sharpening cooperative shock
follows. The two features move \emph{independently} with material parameters, and
their separation at the outlet spans \nGapMin--\nGapMax\ in normalised time --- a
factor of \nGapRange\ across the legacy parameter space, where it shapes the field
every surrogate must reproduce. On the larger dataset that carries most rungs, the
picture is less clean, and we say so here rather than only in
Section~\ref{sec:limitations}: at its stored time resolution the Henry wave crosses
the column between the first two samples, and the second front is mostly a gradual,
dispersive rise rather than a sharp shock. The two-wave structure is therefore the
mechanism behind the legacy design's fields and behind the coordinate change of
Section~\ref{sec:warp}, which is measured there; on the larger dataset it is present
but not resolved well enough to align, and no verdict on that dataset depends on it.

This is established behaviour, not our observation. LeVan and Carta \cite{levan2019adsorption} give the
mechanism: an unfavourable dimensionless isotherm admits a simple wave because low
concentrations travel faster than high ones, a favourable isotherm admits no simple
wave and the solution is a shock, and a favourable isotherm with mass-transfer
resistance approaches a constant pattern. Bozbiyik and co-workers \cite{bozbiyik2017stepped} measured a stepped
breakthrough on aluminium fumarate, and report that it is feed-condition dependent
--- a single step at low water partial pressure, a stepped profile at higher partial
pressure.

\begin{figure}[t]
\centering
\includegraphics[width=\textwidth]{../figures/Fig7_anchor.pdf}
\caption{The reference solver against a published MOF-303 packed bed, with no parameter fitted to the curve. The measured effluent shows an early leak, from uptake on the low-humidity (primary) sites, to an intermediate plateau near \nAnchorPlateauExp{} of the feed, then a cooperative shock much later in the run --- the two-wave structure of Section~\ref{sec:physics}, in a real bed at conditions that are fully on record. The two named discrepancies are visible: the modelled first wave arrives far too early, which the isotherm's low-humidity branch sets, and the shock is too dispersed, which the axial-dispersion closure sets on a bed about \nAnchorPelletsDeep{} pellets deep.}
\label{fig:anchor}
\end{figure}

\subsection{The solver against a real MOF-303 bed, with nothing fitted}
\label{sec:anchor}

Lassitter and co-workers \cite{lassitter2024mass} report a MOF-303 packed-bed breakthrough whose conditions
are fully specified in their supplementary information: a \nAnchorBedLenMm\,\si{\milli\metre}
bed in a \nAnchorTubeMm\,\si{\milli\metre} tube, ambient air at \nAnchorRH\,\% relative
humidity, \nAnchorTK\,\si{\kelvin}, bulk density \nAnchorRhoB\,\si{\kilo\gram\per\metre\cubed}.
The comparison was pre-registered before it was run. Inputs are the cited MOF-303
isotherm \cite{hanikel2021evolution,fathieh2018practical}, their bed, and declared bands for what the source does not give.

The measured curve shows first leak at \nAnchorTfiveExp\ minutes, a Henry-branch
plateau at \nAnchorPlateauExp{} of the inlet, and 50\,\% arrival at
\nAnchorTfiftyExp\ minutes. Our solver, with nothing tuned, gives 50\,\% arrival at
\nAnchorTfiftyModel\ minutes and a plateau at \nAnchorPlateauModel, with
\(\mathrm{nRMSE} = \nAnchorNrmse\) against the digitised effluent points --- against
\nAnchorNrmseComsol\ for the authors' own two-dimensional COMSOL model with fitted kinetics.
The two-wave breakthrough of Section~\ref{sec:physics} is what a real MOF-303 bed
does, and the cited isotherm produces it --- though, as the next paragraph says,
not its first-wave timing (Figure~\ref{fig:anchor}).

Two discrepancies, both attributable, both reported. The first wave arrives
\emph{too early} in the model (\nAnchorTfiveModel\ minutes against
\nAnchorTfiveExp\ measured); nothing in the declared band moves it, and what sets it
is the isotherm's low-relative-humidity branch, whose Henry fraction was fitted to
the step position and the capacity only. This is the first dynamic measurement that
constrains it, and it is a limitation of our MOF-303 parameterisation rather than of
the solver. The shock is \emph{too dispersed}: the axial-dispersion closure gives a
column P\'eclet number of \nAnchorPeCol\ on a bed about \nAnchorPelletsDeep{} pellets deep, where a
packed-bed correlation does not apply. That discrepancy is now measured rather than
attributed --- Section~\ref{sec:limitations} shows that \nAnchorTcutPct\,\% of the
excess arrival time comes from evaluating one correlation at a limit this bed does
not satisfy.

%% ===========================================================================
\section{Dataset, splits and statistics}
\label{sec:dataset}

\begin{figure}[t]
\centering
\includegraphics[width=\textwidth]{../figures/Fig2_dataset.pdf}
\caption{Dataset v2. \textbf{A}~The material space, \emph{one point per material} --- the legend counts are computed from the split and asserted to partition it, because the figure this replaces iterated over samples while plotting material descriptors and showed no training points at all (Section~\ref{sec:corrections}, B27). \textbf{B}~Damk\"ohler per material, the axis rung L6's primary estimand is regressed on. \textbf{C}~The accessible Henry region per material: Henry's law holds for every material, and the region over which it dominates does not (B57).}
\label{fig:dataset}
\end{figure}

\subsection{Design}

The dataset contains \nDataNsamples\ accepted simulations over \nDataNmaterials\
materials and \nDataNconditions\ operating conditions, sampled by a scrambled Sobol
sequence over eight material parameters and three condition parameters (Figure~\ref{fig:dataset});
\nDataNrejected\ candidates were rejected by a screen before being solved. The
kinetic coefficient is not sampled independently: it is derived from particle size
through a Glueckauf relation \cite{glueckauf1947theory,glueckauf1955theory}, so kinetics
and geometry stay physically consistent. The later of those two papers is closed at the
publisher and was not read; the linear-driving-force factor is verified through Perry's
handbook \cite{levan2019adsorption}, and the earlier paper is cited because the plain
factor is attributed to it (Section~\ref{sec:reproducibility}).

Damk\"ohler number is reported with its denominator attached, because the two useful
summaries differ, and quoting either unlabelled is how a reader ends up unable to
reconcile them. \textbf{Per sample} (\nDataNsamples\ values): median
\nDaSampleMedian, range \nDaSampleMin--\nDaSampleMax, \nDaSampleBand\,\% in the
informative band. \textbf{Per material} (\nDataNmaterials\ values, each the median
over that material's own conditions --- and the axis rung L6's slope test regresses
on): median \nDaMatMedian, range \nDaMatMin--\nDaMatMax, spanning \nDaMatDecades{}
decades, \nDaMatBand\,\% in band.

That range is the reason this dataset exists. An earlier dataset had a per-sample
median of \nDaLegacyMedian\ and placed \nDaLegacyAbove\,\% of its samples
\emph{above} the band, which made the primary hypothesis of rung L6 --- that
separating the kinetic object should pay most where kinetics dominate --- untestable
in principle rather than merely underpowered. That figure had itself lived only in a
comment in the generator until it was recomputed for this paper
(Section~\ref{sec:corrections}, B63); it reproduces exactly.

\subsection{The splits, and why the test is clustered}

\paragraph{The statistics in plain terms.} Several breakthrough fields share one
material, so they are not independent observations; we therefore resample whole
materials, never individual fields. We set each significance threshold so that, in
simulations where no real effect exists, the test falsely claims one only rarely, at
the rate stated below for each design. For every ``no difference'' we state the
smallest effect the test could have detected, and where a model was chosen as the
best of several on the test materials we widen its interval to pay for that choice.

Materials are held out, not conditions. Every rung on this dataset reports over the
same five-fold partition of the 240 materials, so that L1, L2, L6 and L7 are
compared over the same held-out set in the same folds.

Errors are strongly clustered by material, and a test that ignores that
over-rejects. Measured under a realistic null, the percentile cluster bootstrap
rejects \nFprLegacyNominal\,\% of the time at a nominal 5\,\% level with
\nNclustLegacy\ clusters, and \nFprFoldsNominal\,\% with \nNclustFolds. We therefore
use a calibrated level: \(\alpha = \nAlphaFolds\) at \nNclustFolds\ clusters, which
achieves an empirical size of \nFprFoldsCal\,\%. One published verdict did not
survive that calibration and was withdrawn (Section~\ref{sec:corrections}, A13).

\paragraph{Three designs, their levels, all stated.} Rungs L1, L2, L6 and L7, the
cost accounting and the re-measurement of L5's basis-versus-map decomposition run on the
\nNclustFolds-material folds at \(\alpha = \nAlphaFolds\). Rung L4 runs on the single
split of \nLfourbMatNmat\ held-out materials at \(\alpha = \nAlphaManifest\). Rung L3,
L5's operator comparisons, the first measurement of its decomposition and the
coordinate change run on the earlier \nNclustLegacy-material design and report at \(\alpha = \nLfiveAlpha\),
which achieves an empirical size of \nFprLegacyUsed\,\% there --- stricter than the
\(\alpha = \nAlphaLegacySel\) the calibration selected for that design. The nominal
levels mislead if compared directly: the smaller design's is a quarter of the larger
design's, but its empirical size is \emph{higher} (\nFprLegacyUsed\,\% against
\nFprFoldsCal\,\%), because the percentile bootstrap over-rejects more at fewer
clusters. We state both because a reader cannot otherwise tell which test produced
which interval.

And we state the cost of the smaller design, because it is the weakest statistical
statement in the paper. At \nNclustLegacy\ clusters the design's minimum detectable
effect at 80\,\% power is \nLegacyMDE\,\%: a ``no difference'' verdict from it
excludes only very large effects, and the calibration script says so in those words.
Where a null is reported on that design below, the interval itself is quoted in
relative terms, because the observed paired interval is much tighter than the
design-level figure --- the arms being compared are trained on the same data and
their errors are correlated, which the generic power simulation cannot know.
Section~\ref{sec:limitations} says what else follows.

\subsection{Every null carries its minimum detectable effect}

A null verdict without a power statement is not a result, and this project learned
it from its own error: an audit reported one of our nulls as well powered, at 96\,\%
power against a 20\,\% effect, and the recomputation with material-level clustering
preserved gave seven per cent (A22). A second correction, in the opposite direction,
found that the corrected power simulation itself inflated within-material noise by
\(\sqrt{2}\) (A25). Both are in Section~\ref{sec:corrections}. Every null reported
below carries a minimum detectable effect computed by a single, self-tested
implementation from the real paired differences.

%% ===========================================================================
\section{Methods}
\label{sec:methods}

Every design constant below is read by \texttt{design\_constants.py} from the
dataset manifests, the generator's own physics builder and the solver's signature,
so that this section cannot drift from the code it describes.

\subsection*{Governing model}

The reference is a one-dimensional, non-isothermal packed bed with axial dispersion,
linear-driving-force kinetics and an energy balance with wall loss. With \(c(z,t)\)
the gas-phase adsorbate concentration, \(q(z,t)\) the loading per unit particle mass,
\(T(z,t)\) the bed temperature, \(v\) the \emph{superficial} velocity and
\(u = v/\varepsilon_t\) the interstitial one, the equations integrated are
\begin{align}
\frac{\partial c}{\partial t} &= D_L \frac{\partial^2 c}{\partial z^2}
 - \frac{v}{\varepsilon_t}\frac{\partial c}{\partial z}
 - \frac{1-\varepsilon_t}{\varepsilon_t}\,\rho_p \frac{\partial q}{\partial t},
\label{eq:mass}\\
\frac{\partial q}{\partial t} &= k_{\mathrm{LDF}}\bigl(q^{*}(c,T) - q\bigr),
\label{eq:ldf}\\
C_{\mathrm{b}}\frac{\partial T}{\partial t} &= k_z \frac{\partial^2 T}{\partial z^2}
 - v\rho_g C_{pg}\frac{\partial T}{\partial z}
 + (1-\varepsilon_t)\rho_p(-\Delta H)\frac{\partial q}{\partial t}
 - \frac{4h_w}{D_{\mathrm{in}}}\,(T - T_w),
\label{eq:energy}
\end{align}
with bed heat capacity \(C_{\mathrm{b}} = \varepsilon_t\rho_g C_{pg} +
(1-\varepsilon_t)\rho_p C_{ps}\) and \(q^{*}\) the dual-site isotherm of
equation~\eqref{eq:isotherm}. Here \(\varepsilon_t\) is the \emph{interparticle} (bed)
void fraction and \(\rho_p\) the pellet density, pores included, so the solid term
counts each pellet once. Gas held inside the pellet pores is not carried as a
separate phase; it holds a negligible amount of water next to the adsorbed phase. The cluster-to-primary affinity ratio is fixed at
\nDcBH, and the cluster affinity is set per run so that the cooperative step sits at
the material's step humidity at the feed temperature. Feed concentration follows
from relative humidity through the ideal-gas law and a Magnus saturation pressure.

\paragraph{Boundary and initial conditions.} Both inlets are of flux (Danckwerts)
type: on the feed side of the inlet face there is no dispersion or conduction, so the
feed's advective flux equals the bed's advective plus dispersive (or conductive) flux
just inside it,
\begin{equation}
\frac{v}{\varepsilon_t}\,c_{\mathrm{in}} = \frac{v}{\varepsilon_t}\,c - D_L\frac{\partial c}{\partial z},
\qquad
v\rho_g C_{pg}\,T_{\mathrm{in}} = v\rho_g C_{pg}\,T - k_z\frac{\partial T}{\partial z}
\qquad (z = 0),
\label{eq:danckwerts}
\end{equation}
with zero gradient at the outlet \(z = L\). The bed starts clean, \(c = q = 0\), at
the wall temperature, which is held at the feed temperature.

\paragraph{Fixed properties.} Every run uses \(L = \nDcL\)\,m, \(D_{\mathrm{in}} =
\nDcDin\)\,m, \(\rho_g = \nDcRhog\)\,kg\,m\(^{-3}\), \(C_{pg} = \nDcCp\) and \(C_{ps} =
\nDcCps\)\,J\,kg\(^{-1}\)\,K\(^{-1}\), \(k_z = \nDcKz\)\,W\,m\(^{-1}\)\,K\(^{-1}\) and
\(h_w = \nDcHw\)\,W\,m\(^{-2}\)\,K\(^{-1}\). Axial dispersion follows the packed-bed
correlation \(D_L = 0.7D_m + 0.5\,d_p v/\varepsilon_t\) with \(D_m =
\nDcDmDisp\times10^{-5}\)\,m\(^2\)\,s\(^{-1}\). In the larger dataset the kinetic
coefficient is not sampled but derived from the particle by the Glueckauf relation,
\begin{equation}
k_{\mathrm{LDF}} = \frac{15\,\varepsilon_p D_m/\tau_p}{(d_p/2)^2\,\rho_p\,K},
\qquad K = \frac{q^{*}(c_{\mathrm{in}},T_{\mathrm{in}})}{c_{\mathrm{in}}},
\label{eq:glueckauf}
\end{equation}
with pellet porosity \(\varepsilon_p = \nDcGlEpsp\), tortuosity \(\tau_p = \nDcGlTau\)
and \(D_m = \nDcGlDm\times10^{-5}\)\,m\(^2\)\,s\(^{-1}\). The generator therefore uses
two slightly different values of the molecular diffusivity, one in each correlation;
we report both rather than reconcile them after the data were generated.

\paragraph{Numerics and storage.} Equations~\eqref{eq:mass}--\eqref{eq:energy} are
discretised by the method of lines on \nGridNz\ cell-centred finite volumes
(first-order upwind advection, central dispersion and conduction) and integrated with
SciPy's implicit BDF method at relative and absolute tolerances
\(10^{\nDcRtolExp}\) and \(10^{\nDcAtolExp}\) to \nDcSnapshots\ output times;
Section~\ref{sec:reference} verifies the solver. Each run is integrated to
\nDcHorizon\ stoichiometric times; if the outlet has not reached the acceptance level
the horizon is extended to \nDcHorizonExt\ stoichiometric times in turn, which
\nDcExtendedPct\,\% of the accepted runs of the larger dataset needed. A draw is
rejected before solving if its feed loads too little or too much of capacity, its
Henry constant is not finite and positive, its isothermal mass-transfer zone is not
resolved by the grid, or its stoichiometric time is implausible; \nDataNrejected\ of
the larger dataset's draws were rejected and \nDataNsamples\ accepted. Fields are
stored as \(c/c_{\mathrm{in}}\), \(q/q_{\max}\) and \(T/T_{\mathrm{in}}\) on a uniform
grid in \(z/L\) and \(t/t_{\mathrm{final}}\), at \(\nDcStoreV\times\nDcStoreV\) for the
larger dataset and \(\nDcStoreLegZ\times\nDcStoreLegT\) for the legacy one.

\paragraph{Damk\"ohler number.} The Damk\"ohler number used throughout is computed per
sample as \(\mathrm{Da} = k_{\mathrm{LDF}}\,t_{\mathrm{final}}\), and rung L6 regresses
on its per-material median. Because the horizon is a multiple of the stoichiometric
time that depends on whether it was extended, Da carries that multiple as well as the
kinetic ratio \(k_{\mathrm{LDF}}t_{\mathrm{st}}\); for the \nDcExtendedPct\,\% of runs
that were extended it is larger than the kinetic ratio by the extension factor. We keep the
definition the pre-registration used and state this property rather than redefine
the axis after the result.

\subsection*{Parameter space}

Materials and operating conditions are sampled separately and crossed,
\nNclustFolds\ materials at \nDcNcond\ conditions each, so that a material is a real
grouping in the data (Table~\ref{tab:params}). Each set is a scrambled Sobol design
mapped linearly onto its ranges, with seeds \nDcSeedMat\ (materials) and \nDcSeedCond\
(conditions). The legacy dataset used by rungs L3 and L5 and by the coordinate change
differs in three ways: independent uniform draws, \(k_{\mathrm{LDF}}\) sampled
directly on \([\nDcKldfLo, \nDcKldfHi]\)\,s\(^{-1}\) instead of being derived from
the particle, and a fixed particle size; it has \nDcNmatLeg\ materials and
\nDcNacceptLeg\ accepted runs.

\paragraph{Where real water-harvesting frameworks sit.} The sampled space is meant to
contain the adsorbents a water harvester would use, so we placed published ones in it
(Figure~\ref{fig:watermofs}), on the four axes that are properties of the framework
rather than of the pellet and the packing: step position, capacity, heat of
adsorption and the fraction of capacity taken below the step. The values come from
the primary literature, MOF-303 \cite{hanikel2019rapid,fathieh2018practical},
MOF-801 \cite{furukawa2014water}, CAU-10-H \cite{solovyeva2021cau,frohlich2016water},
aluminium fumarate \cite{hanikel2019rapid,laurenz2020novel}, MIL-160(Al)
\cite{solovyeva2021mil,aumond2026adsorption,lenzen2019metal} and
Co\(_2\)Cl\(_2\)(BTDD) \cite{rieth2017record}; each point spans the values its
sources report, which differ in sample, temperature and in how the step is defined,
and the below-step fractions are mostly read from published figures and are rough.
Of the \nWmofN, \nWmofInside\ lie inside the sampled ranges on every axis. The
exceptions are at the edges, and one is far outside: MOF-801's step, at about
\nWmofEightStep, sits just below the sampled floor of \nWmofStepLo; one reported heat
for CAU-10-H, \nWmofCauHeat\,kJ\,mol\(^{-1}\), lies just beyond the sampled range; and
Co\(_2\)Cl\(_2\)(BTDD), at about \nWmofBtddQ\,mol\,kg\(^{-1}\), holds
\(\nWmofBtddRatio\times\) the largest capacity we sampled, \nWmofQmaxHi. The verdicts of this paper are therefore
statements about frameworks like MOF-303, aluminium fumarate and MIL-160, and should
not be carried to record-capacity materials without new data. The cooperativity
exponent \(n\), the steepness of the step and arguably the property a harvester
designer cares about most, is not placed: the sources report isotherms, not fitted
step exponents, and fitting equation~\eqref{eq:isotherm} to digitised curves would be
a new measurement, not a reading of the literature.

\begin{figure}[t]
\centering
\includegraphics[width=\textwidth]{../figures/FigS2_water_mofs.pdf}
\caption{Published water-harvesting frameworks against the sampled parameter space
(shaded boxes). \textbf{A}~Step position and capacity. \textbf{B}~Fraction of capacity
taken below the step and heat of adsorption. Bars span the values reported across
sources; values, locations in each paper and confidence are recorded in
\texttt{research/water\_mof\_parameters.json}.}
\label{fig:watermofs}
\end{figure}

\begin{table}[t]
\centering
\caption{The parameter space of the larger dataset. A ``material'' here is the upper
block: an idealised framework pellet described by its isotherm (capacity, step
position, step steepness, the fraction of capacity on the low-humidity sites, heat)
and by pellet size, density and packing. Capacity \nDcQmaxLo--\nDcQmaxHi\,mol\,kg\(^{-1}\)
is \nDcQmaxGgLo--\nDcQmaxGgHi\,g of water per g. The step is re-anchored at each feed
temperature so that it sits at the stated relative humidity, so one material's
isotherm is not the same curve at every condition. All eleven parameters are sampled by
scrambled Sobol and mapped linearly; every surrogate except rung L6 receives exactly
these values, standardised on its own fitting set, and L6 receives observable
descriptors instead.}
\label{tab:params}
\begin{tabular}{llll}
\toprule
parameter & symbol & units & range \\
\midrule
\multicolumn{4}{l}{\emph{material}} \\
saturation capacity & \(q_{\max}\) & mol\,kg\(^{-1}\) & \nDcQmaxLo--\nDcQmaxHi \\
heat of adsorption & \(\Delta H\) & kJ\,mol\(^{-1}\) & \nDcDHLo\ to \nDcDHHi \\
step position (relative humidity) & \(\mathrm{RH}_{\mathrm{step}}\) & -- & \nDcStepLo--\nDcStepHi \\
cooperativity exponent & \(n\) & -- & \nDcNexpLo--\nDcNexpHi \\
Henry-site fraction & \(f_H\) & -- & \nDcFhLo--\nDcFhHi \\
particle diameter & \(d_p\) & mm & \nDcDpLo--\nDcDpHi \\
particle density & \(\rho_p\) & kg\,m\(^{-3}\) & \nDcRhopLo--\nDcRhopHi \\
bed porosity & \(\varepsilon_t\) & -- & \nDcEpsLo--\nDcEpsHi \\
\midrule
\multicolumn{4}{l}{\emph{operating condition}} \\
feed relative humidity & \(\mathrm{RH}_{\mathrm{feed}}\) & -- & \nDcRHLo--\nDcRHHi \\
superficial velocity & \(v\) & m\,s\(^{-1}\) & \nDcVLo--\nDcVHi \\
feed temperature & \(T_{\mathrm{in}}\) & K & \nDcTinLo--\nDcTinHi \\
\bottomrule
\end{tabular}
\end{table}

\subsection*{Error metric and statistical test}

The primary metric is a per-sample normalised root-mean-square error on the
concentration channel. For sample \(s\), with \(\hat c_{ij}\) and \(c_{ij}\) the
predicted and reference \(c/c_{\mathrm{in}}\) on the rung's evaluation grid,
\begin{equation}
\mathrm{nRMSE}_c(s) = \frac{\Bigl[\frac{1}{N_zN_t}\sum_{i,j}(\hat c_{ij}-c_{ij})^2\Bigr]^{1/2}}
{\max_{i,j} c_{ij} - \min_{i,j} c_{ij}} ,
\label{eq:nrmse}
\end{equation}
the RMSE over the whole space--time field divided by that sample's own reference
range. Every accepted run starts from a clean bed and ends with the outlet near the
feed concentration, so the denominator is close to one and cannot vanish. Where only
part of the field is scored (the time axis of rung L4) the numerator runs over the
scored window and the denominator stays the range of the full field. The same quantity
is computed for \(q/q_{\max}\) and \(T/T_{\mathrm{in}}\) and reported as secondary; the
channels are never pooled. Rung L7 scores the outlet curve only. Per-sample errors are
averaged over the training seeds, pooled across held-out folds so that each sample
appears once, and averaged with equal weight per sample.

Two arms are compared on identical samples by a paired, cluster-robust percentile
bootstrap: materials are resampled with replacement, all conditions of each drawn
material are kept, and the mean paired difference is recomputed. A difference is
reported as significant only if the percentile interval excludes zero, at a level
calibrated per design: simulating clustered nulls with each design's real
conditions-per-material counts, we take the largest nominal level on a fixed grid
whose empirical false-positive rate does not exceed the target. This gives
\(\alpha = \nAlphaFolds\) for the five-fold design (empirical size \nFprFoldsCal\,\%);
on the legacy design we report at the stricter \(\alpha = \nLfiveAlpha\) (empirical
size \nFprLegacyUsed\,\%). The time axis of rung L4 has an intermediate number of
clusters, at which no calibration was run; it uses the level both neighbouring
designs select. Rung L6's primary estimand is the least-squares slope of the
per-material paired difference on \(\log_{10}\mathrm{Da}\), with a material bootstrap
interval. Best-of-grid selection is discussed in Section~\ref{sec:limitations}.

\subsection*{Surrogates, encodings and evaluation designs}

Three evaluation designs are used, and each result states which. In the
\emph{five-fold} design the \nNclustFolds\ materials of the larger dataset are
partitioned into five folds, every material is held out once, and each fit sees the
other four folds. The \emph{single split} holds out the dataset's own
\nLfourbMatNmat\ test materials. The \emph{legacy} design holds out \nNclustLegacy\
of the earlier dataset's \nDcNmatLeg\ materials. Every learned arm is trained with
three seeds and fitted --- including any proper-orthogonal-decomposition basis and any
scaler --- on its fitting set only, which is asserted disjoint from the held-out
materials before every fit.

\textbf{Absolute errors are not comparable across rungs.} The rungs differ in
dataset, in field encoding, in what is scored (the whole field, a time window or the
outlet curve) and in held-out set. Only paired comparisons within a rung are
inferential, and Table~\ref{tab:rungs} states each rung's encoding so that no
cross-rung reading is made by accident. The one exception is stated where it is used
(Section~\ref{sec:cost}): L5's operators and the coordinate change share a split, an
encoding and a metric, so their absolute errors are comparable.

\begin{table}[p]
\centering
\caption{Surrogates by rung. Absolute errors in different rows are not comparable.}
\label{tab:rungs}
\footnotesize
\renewcommand{\arraystretch}{1.15}
\resizebox{\textwidth}{!}{%
\begin{tabular}{@{}p{0.07\textwidth}p{0.13\textwidth}p{0.25\textwidth}p{0.30\textwidth}p{0.20\textwidth}@{}}
\toprule
rung & data, design & input \(\to\) output; encoding & architecture & training \\
\midrule
L1 & larger, five-fold &
parameters \(\to\) POD coefficients per channel; field subsampled from the stored grid &
ridge; random forest; gradient boosting; perceptron with standardised targets &
scikit-learn; perceptron Adam with early stopping \\
L1 curve & larger, five-fold & as L1, training materials subsampled & the L1 perceptron & as L1 \\
L2 & larger, five-fold & as L1 &
perceptron width and depth, boosting depth, forest leaf size, swept in both directions & as L1 \\
L3 & legacy &
parameters \(\to\) POD coefficients; full stored field &
perceptron and two Kolmogorov--Arnold families (Gaussian radial basis, Chebyshev), depth \nMthThreeDepth, matched to \nMthThreeBudgets\ parameters; widths \nMthThreeWidths &
Adam, cosine schedule, \nMthThreeSteps\ steps, per-family learning-rate grid \\
L4 & larger; time axis on the training materials, material axis on the single split &
(parameters, \(z\), \(t\)) \(\to\) \((c,q,T)\) pointwise, bounded output head &
one perceptron, \nMthFourDepth\ layers of \nMthFourWidth\ (\nMthFourParams\ parameters); data-only twin and physics-weighted arms (fixed weights, gradient-norm balancing, NTK and self-adaptive weighting); the residual of equations \eqref{eq:mass}--\eqref{eq:energy} with \eqref{eq:danckwerts}, as corrected under B72 &
Adam, cosine schedule, \nMthFourSteps\ steps of \nMthFourSup\ data, \nMthFourColloc\ collocation and \nMthFourBc\ boundary points; then L-BFGS (polish) \\
L5 & legacy &
parameters \(\to\) \(c\) field, subsampled from the stored field &
DeepONet and DeepOKAN over basis size \(p\); Fourier neural operator over mode count; all matched to \nMthFiveBudget\ parameters &
Adam, cosine schedule, \nMthFiveSteps\ steps, per-arm learning-rate grid \\
L5 basis vs.\ map & legacy, then re-measured on larger, five-fold &
parameters \(\to\) POD coefficients of \(c\) over basis size &
one gradient-boosted regressor per mode & three seeds \\
L6 & larger, five-fold &
\nMthSixNdesc\ observable descriptors (isotherm at three temperatures, bed and operating values), \(z\), \(t\) \(\to\) \((c,q,T)\) &
joint perceptron versus separate equilibrium, kinetic and transport heads, with \(q\) from equation~\eqref{eq:ldf}; both \nMthSixBudget\ parameters, depth \nMthSixDepth &
Adam, cosine schedule, \nMthSixSteps\ steps \\
L7 & larger, all samples; paired with L1 out of fold &
parameters \(\to\) outlet curve &
equilibrium shock, Klinkenberg, constant pattern; nothing fitted & none \\
Warp & legacy &
parameters \(\to\) \(c\), in a fixed frame or the two-front frame of equation~\eqref{eq:warp} &
POD plus one gradient-boosted regressor per mode; front trajectories from a POD of the landmark curves &
three seeds; oracle arm given the true fronts \\
\bottomrule
\end{tabular}}
\end{table}

%% ===========================================================================
\section{The ladder}
\label{sec:ladder}

\begin{figure}[t]
\centering
\includegraphics[width=\textwidth]{../figures/Fig1_ladder.pdf}
\caption{The falsification ladder. \textbf{A}~Each rung's intervention as a
percentage of that rung's own control, with its cluster-robust confidence interval;
positive means the intervention reduces held-out error. Where an arm was chosen as
the best of a grid on the held-out materials (L2, L5's operator comparison), the
interval drawn is the one that pays for that choice; L3 draws its strongest single
pair per pair, and its note gives the count that survives selection. Dashed ticks
mark a minimum detectable effect where the rung's analysis reports one; for each null
the text states what its interval excludes. The physics-as-a-loss rung is withdrawn (B72) and draws no
bar. The oracle arm is drawn hollow because it is handed the answer and is an upper
bound, not a comparable result. \textbf{B}~The primary
pre-registered estimand of rung L6, in its own units. \textbf{C}~Verification of the
reference (Table~\ref{tab:l0}).}
\label{fig:ladder}
\end{figure}

\begin{figure}[t]
\centering
\includegraphics[width=\textwidth]{../figures/Fig3_rungs.pdf}
\caption{The rungs that measure a curve. \textbf{A}~L1's fixed-basis arms over 240 held-out materials, against the proper-orthogonal-decomposition floor. \textbf{B}~The materials learning curve with the basis refit on each subset and with it fixed on all 192; the two lie on top of each other. \textbf{C}~L2's capacity sweep: held-out error is U-shaped in depth while training error falls monotonically. \textbf{D}~L3 at three matched budgets, each family at its own best learning rate --- the in-bar labels are those rates, and they span three orders of magnitude.}
\label{fig:rungs}
\end{figure}

\subsection{L1 --- more data}

\paragraph{The arms.} Pooled out of fold over 240 held-out materials, with 64
proper-orthogonal-decomposition modes per channel, the multilayer perceptron is the
best fixed-basis arm at \nLoneMLP, ahead of gradient boosting at \nLoneXGB, random
forest at \nLoneRF\ and ridge at \nLoneRidge. Against gradient boosting the
difference is \nLoneMLPvsXGBdiff\ with confidence interval
[\nLoneMLPvsXGBlo, \nLoneMLPvsXGBhi] --- significant, and the perceptron is ahead in
every one of the five folds. This reverses the earlier result at 48
materials, where the three nonlinear arms were indistinguishable and the perceptron
had the worst point estimate.

The floor moved too. The proper-orthogonal-decomposition floor is \nLonePodFloor,
and the best arm sits \(\nLoneXfloor\times\) above it. On the smaller dataset the same comparison gave \nLoneLegacyXfloorMin--\nLoneLegacyXfloorMax\(\times\), against \(\nLoneXfloor\times\) now, so \emph{the basis is even less the constraint than it was}: 192 training materials instead of 48 improved the arm, and improved the floor beneath it by more.

\paragraph{The learning curve, and the rung's actual verdict.} Held-out error (Figure~\ref{fig:rungs}B) falls
\nLCsmallest\ \(\to\) \nLClargest\ as training materials go 12 \(\to\) 192, a factor
of \nLCratio. Every consecutive step is a significant improvement \emph{including
the last}: 96 \(\to\) 192 gives \nLClastDiff\ with interval
[\nLClastLo, \nLClastHi]. In the number of training materials \(N_{\mathrm{mat}}\) the
error falls as \(N_{\mathrm{mat}}^{-\nLCbeta}\), exponent interval
[\nLCbetaLo, \nLCbetaHi]. Extrapolated at that exponent --- \emph{outside} the range
we claim, and stated as an extrapolation because the next sentence forbids it ---
halving the error would cost about \(\nLCmaterialsForHalving\times\) more materials.

\textbf{Verdict, in the pre-declared words: the materials axis is NOT ELIMINATED at
the largest size available.} We claim no extrapolation beyond twice the measured
range.

\paragraph{What more materials buy, and what they do not.} Refitting the basis on
each subset changes nothing: \nLCsmallestFixed\ against \nLCsmallest\ at 12
materials, \nLClargestFixed\ against \nLClargest\ at 192. A basis fitted on 12
materials reconstructs held-out fields as well as one fitted on 192, once the
regressor sees the same data. \emph{None of the gain is in the basis; all of it is
in the parameter-to-coefficient map.} Modes whose coefficient map clears
a fifth of the variance grow from \nLCmodesSmall\ to \nLCmodesLarge\ and are still climbing.

The conditions axis behaves differently: exponent \nLCcondBeta, \nLCexpRatio{} of the material
exponent. Both axes are open; materials are the one that matters.

\subsection{L2 --- more capacity}

\nLtwoNconfigs{} configurations in both directions around the L1 setting, on the same
folds and basis. \nLtwoNUshaped{} families are U-shaped with interior optima
(width \nLtwoWidthBest{} for the perceptron's width, \nLtwoDepthBest{} layers for its depth, tree
depth \nLtwoXgbBest{} for gradient boosting) and the random forest
\nLtwoRfShape\ --- it is still improving at its last step, but that step is a fully
grown forest, the ceiling of its capacity knob. The decisive evidence is decoupling:
gradient boosting drops its training error by more than an order of magnitude across
the depth sweep while held-out error passes through a minimum and then rises.

\textbf{Verdict, in the pre-declared words: capacity is NOT ELIMINATED.} The rule
eliminates it only if every family has saturated, and one has not: the random forest
is still improving at its last step (fully grown trees, \nLtwoRfBest; last-step
interval [\nLtwoRfLastLo, \nLtwoRfLastHi]). The rule then asks for the sweep to be
extended, and it cannot be --- one sample per leaf is the ceiling of that family's
capacity --- so the verdict stays as declared. Our reading, labelled as a reading and
not as the verdict: the \nLtwoNUshaped{} neural and boosting families saturate with interior
optima, and the one family that has not run out of improvement has run out of
capacity while sitting well behind the best perceptron. \textbf{Either way, the optimum is not where the
earlier sweep left it.} The best configuration anywhere is \nLtwoBestFam/\nLtwoBestCfg\
at \nLtwoBest, which is \(\nLtwoVsLonePct\,\%\) better than the best L1 arm, the
perceptron (interval \(\nLtwoVsLoneLo\)--\(\nLtwoVsLoneHi\,\%\); adjusted for having
picked the best of \nLtwoSelK\ configurations on the same materials --- the
\nLtwoNconfigs\ of the sweep less the L1 setting, which appears in it twice and is the
reference ---
\(\nLtwoSelLo\)--\(\nLtwoSelHi\,\%\), Section~\ref{sec:limitations}). At 48 materials the
same comparison --- best configuration against the best L1 arm, there gradient
boosting --- gave \(\nLtwoLegacyVsLonePct\,\%\), with six and eight layers in a
near-tie (\nLtwoLegacyDsix{} against \nLtwoLegacyDeight) where at 192 the optimum is
eight. The earlier sentence ``capacity is irrelevant'' was scoped to a material count
without saying so, and is withdrawn as A26. The wall stands: the best fixed-basis
surrogate still sits \(\nLtwoXfloor\times\) above the floor.

What changes with the material count is \emph{which} configuration is best, not that
capacity grows: the perceptron's depth optimum moves up from a near-tie to eight
layers, its width optimum moves \emph{down}, and the boosting depth moves by one step,
within a near-tie. We therefore draw no data-limited inference from L2 alone; that
diagnosis rests on the learning curve of L1.

\subsection{L3 --- a better basis}

At matched parameter counts and each family at its own best learning rate, the
multilayer perceptron beats both Kolmogorov--Arnold families at every budget:
\nLthreeSigPair\ of the \nLthreePairs\ pairwise comparisons are significant as pairs. Both arms
of each pair are chosen on the held-out materials, so we also pay for those choices
with one simultaneous interval over all \nLthreeSelK\ pairs of learning rates
(Section~\ref{sec:limitations}). \nLthreeSigSim\ of the \nLthreePairs\ remain
significant, and in none is the Kolmogorov--Arnold network the better arm.
\nLthreeArms\ arms were trained across
\nLthreeNrates\ learning rates spanning \nLthreeDecades\ decades, until every
Kolmogorov--Arnold optimum was bracketed on both sides. Our protocol's rule~4 says an
optimum on the edge of a sweep is admissible only if the metric has saturated there,
and otherwise the sweep is extended. The perceptron selects the lowest rate at every
budget, which is an edge, and it is not saturated in that sense: the next rate up is
worse by \nLthreeEdgeSatMin--\nLthreeEdgeSatMax\,\%, so a still lower rate might do
better. For this verdict the open edge is conservative --- a better-tuned perceptron
could only widen its lead over arms that are all bracketed --- so we report it rather
than extend it. The Kolmogorov--Arnold arms, the ones being argued against, are the
ones whose tuning is complete.

Comparing against the strongest configuration found anywhere rather than the
matched-budget cell --- a separate audit configuration (a fixed budget, depth three,
each arm at the better of two higher learning rates per seed), which is why its
perceptron rate differs from the sweep's: the best Kolmogorov--Arnold network in the entire study is the
\emph{least} Kolmogorov--Arnold-like one, and it reaches \nLthreeBestKan\ against
the perceptron's \nLthreeBestKanMlp\ --- a difference of \nLthreeBestKanDiff,
interval [\nLthreeBestKanLo, \nLthreeBestKanHi] --- \nLthreeBestKanSig. That interval
is per pair: the learning rate of each arm was chosen on the held-out materials, so
it is the count above, not this pair, that carries the selection.

\paragraph{The tuning envelope is itself a result.} The families optimise
\nLthreeLrDecades{} orders of magnitude apart: the Chebyshev family at \nLthreeChebyLr\ and the
perceptron at \nLthreeMlpLr. \nLthreeNfailAll{} configurations failed on every seed and
\nLthreeNfailSome{} more on some seeds; all are excluded and listed, never averaged in. Every family, the perceptron included, was tuned over the same
learning-rate grid; none was left at a default.

\paragraph{The literature is split, and this result is worth more because of it.}
Shukla and co-workers \cite{shukla2024comprehensive} report that spline-based
Kolmogorov--Arnold networks lack
accuracy and efficiency while modified ones on low-order orthogonal polynomials are
\emph{comparable} to perceptron baselines but not robust. But Wang and co-workers'
KINN \cite{wang2025kinn} reports that it ``significantly outperforms MLP regarding accuracy and
convergence speed'' on solid-mechanics problems, and Rigas and co-workers \cite{rigas2026training} report
that residual-gated adaptive Kolmogorov--Arnold networks ``consistently outperform
parameter-matched cPIKANs and PirateNets'' --- and PirateNet is a perceptron
architecture. Kiyani and co-workers \cite{kiyani2025optimizing} find that optimiser
choice dominates accuracy for physics-informed models, Kolmogorov--Arnold networks
included, and that both families improve by orders of magnitude under the \emph{same}
quasi-Newton schemes --- which is why every competitor optimum here is bracketed
rather than left at a default. Nor is a Kolmogorov--Arnold operator new: DeepOKAN
\cite{abueidda2025deepokan} is published, data-driven, with the same Gaussian
radial basis, and our claim for that arm is the application to column-scale transfer
and nothing more (A12). An earlier draft of this section cited two of these papers for a
consensus that does not exist (Section~\ref{sec:corrections}, B46). The honest
statement is narrower and more useful: on \emph{this} problem, at matched
parameters, with every competitor optimum bracketed, the perceptron wins all
\nLthreePairs\ comparisons by point estimate and \nLthreeSigSim\ significantly once both
learning-rate choices are paid for; and the
papers that disagree are cited so a reader sees the disagreement rather than
discovers it.

\paragraph{Interpretation.} Not ``Kolmogorov--Arnold networks are bad''. At a fixed
budget, parameters spent on per-edge basis richness buy training fit and not
generalisation, while the same parameters spent on width buy both. The
Kolmogorov--Arnold arms fit the training set as well as or better than the
perceptron at every budget and are worse on held-out materials every time.

\subsection{L4 --- the residual as a loss, and physics at inference}

\paragraph{Read this first: the residual of this rung was the wrong equation.} While
writing the methods of this paper we derived the residual from the code and found
that its gas mass balance advected the gas at the superficial velocity \(v\) instead
of the interstitial \(v/\varepsilon_t\) that the solver integrates, with the same
omission in its inlet condition (Section~\ref{sec:corrections}, B72). Since the gas
holds almost nothing compared with the solid, the front speed is set by that term, so
the residual described fronts moving \(\varepsilon_t\) times too slowly. The test is
direct: evaluated on the solver's own stored fields, the residual this rung used leaves
\nBsevtwoOldResid\,\% of the adsorption term unexplained, against
\nBsevtwoNewResid\,\% for the solver's equation, worse in \nBsevtwoNworse\ of
\nBsevtwoNsamples\ samples. \textbf{Every result below is therefore a measurement of a
wrong equation used as a prior, not of this column's physics, and none of it may be
read as ``physics helps'' or ``physics hurts'' (A28).} We report the results because
they were pre-registered and because they are informative about what a mis-specified
residual does. The corrected residual is now checked against the solver by a gate,
and the rung is being re-run with it under a dated amendment to its pre-registration.

\paragraph{Design.} Ten weighting schemes were pre-registered and compared against a
data-only twin at matched architecture, matched optimiser steps and matched
supervised budget --- fixed weights over four decades, gradient-norm balancing at
three target ratios \cite{wang2021understanding}, neural-tangent-kernel weighting
\cite{wang2022when} (with per-point gradient norms as a trace estimate, a stated
approximation to their eigen-decomposition), and self-adaptive per-point weights
\cite{mcclenny2023self} --- because the earlier version of this rung rested on a single
weighting rule, which is not a defensible basis for the claim that physics hurts. The
pre-registered edge rule (rule~4: an optimum on the edge of a sweep is admissible only
if the metric has saturated there) then extended the fixed-weight sweep one decade
below its lowest weight, and on the material axis a second decade. A physics arm whose
seen-window error exceeds three times the twin's has abandoned the data fit (the B16
signature) and is excluded as failed-to-train, never averaged in: \nLfourbNabandoned\
arms did, leaving \nLfourbNeligible\ eligible on the time axis and, with the second
extension, \nLfourbNeligibleMat\ on the material axis.

\paragraph{Time axis, with the residual used: no difference.} In the
pre-declared words: there is no difference at the best weighting found. The best
physics arm, the lowest fixed weight, reaches \nLfourbTimeBest\ against the twin's
\nLfourbTimeTwin, a paired difference with interval [\nLfourbTimeLo, \nLfourbTimeHi]
over \nLfourbTimeNmat\ materials; the design detects a \nLfourbTimeMDE\,\% effect with
\(80\,\%\) power. The legacy statement that physics as a loss hurts is withdrawn. This
verdict exists only because the edge rule was applied: at the lowest weight of the
original sweep, physics was significantly worse than the twin, and the one further
decade the rule demanded tied it (Section~\ref{sec:corrections}, B69). On the residual
then in use, rung L4 is the clearest case in this paper of a pre-registered rule
changing a result that a reading
of the rule after seeing the data could have declined to test.

\paragraph{Time axis, after polishing: the ordering flips.} The pre-registration's
fourth question polishes both arms with L-BFGS after Adam
\cite{rathore2024challenges}. Both improve, and the twin improves more: after
polishing the twin reaches \nLfourbTimePolTwin\ and the physics arm
\nLfourbTimePolBest, interval [\nLfourbTimePolLo, \nLfourbTimePolHi] --- the data-only
arm is significantly better. The pre-registration declared a flip to be the reported
result, and we report the numbers; with the residual this rung used, they carry no
reading about physics.

\paragraph{Material axis, with the residual used: the residual arm is better.} In the
pre-declared words the analyser reports its positive verdict on the material axis; with
the residual this rung used, that sentence means only that this particular prior
improved the fit. The best arm is the second-smallest fixed weight, and it is now an
interior optimum: the decade below it, which the edge rule demanded a second time,
reaches \nLfourbMatBelow\ and the decade above \nLfourbMatAbove, both worse than its
\nLfourbMatBest. Against the twin's \nLfourbMatTwin\ the paired difference (twin minus
physics) has interval [\nLfourbMatLo, \nLfourbMatHi] over \nLfourbMatNmat\ held-out
materials, which excludes no difference. The ordering survives polishing, but on this
axis polishing makes \emph{both} arms worse: after L-BFGS the twin reaches
\nLfourbMatPolTwin\ and the physics arm \nLfourbMatPolBest, interval
[\nLfourbMatPolLo, \nLfourbMatPolHi]. The effect is small --- \nLfourbMatGainPct\,\%
of the error --- and it is confined to a single weight: a decade either side, the
advantage is gone.

\textbf{This verdict carries a selection caveat.} The pre-registration chose the best physics arm on the same held-out
materials it is then tested on --- the generous direction for a null, and the
anti-conservative one for a positive. An isolated best weight among
\nLfourbNeligibleMat\ eligible arms is exactly the shape a selection artefact takes.
We therefore recomputed the interval so that it holds for every eligible arm at once,
and so for the selected one however it was selected (a studentised max-statistic
cluster bootstrap, Section~\ref{sec:limitations}). It widens to
[\nLfourbMatSelLo, \nLfourbMatSelHi] and still excludes no difference. That is a
statement about the residual this rung used and about the procedure, which the
corrected re-run will be put through unchanged; it is not a statement about what the
column's physics does as a prior.

\paragraph{Physics at inference.} Test-time refinement --- minimising the residual for
each held-out unit with no data --- was pre-registered with an invalidation
condition: on the time axis it is anchored to the frozen model's own prediction on the
seen window, and a seen-window error rising by more than \(10\,\%\) means the anchor
is insufficient and the result on that axis is scoped. It rose by
\nLfourbTimeSeenRise\,\% on average. The time-axis refinement result is therefore
scoped, and the material axis --- which has no anchor by design, because nothing about
an unseen material is known --- carries the reportable refinement result. \textbf{With the residual used, refinement
makes transfer to unseen materials worse} --- which is what minimising the wrong
equation should do, and so no evidence about the right one.
Refining the twin with the residual alone takes its held-out error from
\nLfourbRefTwinBefore\ to \nLfourbRefTwinAfter, about \nLfourbRefTwinFactor{} times
worse (paired difference interval [\nLfourbRefTwinLo, \nLfourbRefTwinHi]), and the
best physics-trained arm fares no better, \nLfourbRefBestBefore\ to
\nLfourbRefBestAfter. A labelled post-hoc sweep over refinement scope, learning rate
and step count (Section~\ref{sec:corrections}, B60) asks whether this is a property of
one configuration: over \nLfourbRefSweepN\ configurations, the best change found at
any checkpoint is \nLfourbRefSweepBest\,\% of the held-out error (negative means the
error fell), chosen on the held-out materials and so
the most generous reading available, and \nLfourbRefSweepNworse\ of the
\nLfourbRefSweepN\ are worse than the unrefined model at their last checkpoint.

\begin{figure}[t]
\centering
\includegraphics[width=\textwidth]{../figures/Fig4_operators.pdf}
\caption{The operator rung. \textbf{A}~Held-out error against basis size, with the proper-orthogonal-decomposition floor. The floor falls by more than an order of magnitude and no arm follows it; the Fourier operator has no basis size and is drawn as a horizontal line. \textbf{B}~What ``flat'' means: the paired differences against \(p = 8\), every interval spanning zero across a sixteenfold change in basis size. \textbf{C}~Per-mode \(R^2\) of the parameter-to-coefficient map on held-out materials at the largest basis --- \nLfiveCountModes\ modes clear \(R^2 > 0.2\), the leading \nLfiveLeadModes\ and \nLfiveIsolatedModes\ isolated later ones, and the median \(R^2\) of the \nLfiveTailModes\ modes after the leading run is \(\nLfiveTailMedian\), worse than predicting the mean.}
\label{fig:operators}
\end{figure}

\subsection{L5 --- does an operator remove the wall?}

At matched parameter counts and each family at its own best learning rate over
\nLfiveNrates\ rates spanning \nLfiveDecades\ decades, DeepONet is \textbf{flat in basis size} (Figure~\ref{fig:operators}): \nLfiveFlatDiff\ from
\(p = 8\) to \(p = 128\), interval [\nLfiveFlatLo, \nLfiveFlatHi], which spans zero.
Stated as what it excludes, which is the honest form for a null: a sixteenfold
increase in basis size cannot have improved this arm by more than
\nLfiveFlatImproveMax\,\% or degraded it by more than \nLfiveFlatDegradeMax\,\% of
its own error. That is the bound we claim, and it is not the same as ``no effect''
--- while its own proper-orthogonal-decomposition floor falls from
\nLfiveFloorSmall\ to \nLfiveFloorLarge, a factor of \(\nFloorDrop\). The best
DeepONet found anywhere, at the smallest basis, sits \(\nONetxFloor\times\) above the
floor of the largest. Rule~4 is the rule this rung was most exposed on: at the two
smallest bases DeepONet's best learning rate first sat on the top of the grid, still
moving. The grid was extended one rate upward for those arms, the optimum stayed where
it was and is now interior, and every selected rate in this rung is bracketed --- so
the Fourier operator below is not compared with an under-tuned DeepONet. \emph{The best DeepONet found anywhere, at the smallest basis, is as
accurate as an oracle handed the true coefficients of its first
\nLfiveOracleONet\ modes, and adding basis functions does not move it.}

\paragraph{The arm the theory names.} Lanthaler and co-workers \cite{lanthaler2023nonlinear} prove that operator
architectures with a linear reconstruction step --- DeepONet, PCA-Net, and
proper-orthogonal-decomposition-plus-regressor --- are lower-bounded on
advection-dominated problems, and that the Fourier neural operator escapes that
bound through nonlinear reconstruction. Our original rung tested only the family the
theory condemns; that deviation from our own frozen protocol is recorded as B33. The
prediction was pre-registered before the run.

The Fourier neural operator is more accurate, \nLfiveFNO\ against DeepONet's
\nLfiveONet, but \textbf{not significantly so once the choice is paid for}. Taken as a
pair fixed in advance, the difference of \nLfiveONetFNOdiff\ has interval
[\nLfiveONetFNOlo, \nLfiveONetFNOhi], which excludes zero by a small fraction of
either error. But both arms are the best of grids chosen on the same
\nNclustLegacy\ held-out materials: the best of \nLfiveSelNfno\ Fourier cells against
the best of \nLfiveSelNonet\ DeepONet cells. Paying for both choices at once --- a
studentised max-statistic bootstrap over all \nLfiveSelK\ pairs
(Section~\ref{sec:limitations}), critical value \nLfiveSelQ\ standard errors --- widens
the interval for the selected pair to [\nLfiveSelLo, \nLfiveSelHi], which includes
zero. We reported this comparison as significant in an earlier draft and withdraw
that word (Section~\ref{sec:corrections}, B75). What survives is the direction the
theory predicts, measured on this physical problem rather than in an approximation
bound. \textbf{And it does not remove the wall}: the Fourier operator sits
\(\nFNOxFloor\times\) above the floor where DeepONet sits \(\nONetxFloor\times\). A
\(\nFNOgain\times\) improvement, against a reference level \(\nFNOxFloor\times\) below
it, is a dent and not an escape. That level is the linear-reconstruction floor at
\(p = 128\), and the Fourier operator is not bound by it.

One row deserves its own sentence. At a matched parameter budget the best Fourier
operator spends it on spectral weights rather than on channels: it is
\emph{\nLfiveFNOwidth\ channels wide} against the best DeepONet's
\nLfiveONetWidth, a factor of \nLfiveWidthRatio. A nonlinear reconstruction matches
or beats the best linear one with a fraction of the channel count, at the same
parameter count. It does not follow that reconstruction is what separates the
families in general. Proper orthogonal decomposition plus a pointwise regressor and
DeepONet both reconstruct linearly, and they differ by \(\nONetLever\times\)
(Section~\ref{sec:mechanism}). Most of the learner lever is therefore training the
basis and the coefficients together, not making the reconstruction nonlinear. Note
that this is the \emph{matched-budget} comparison and not a reading of the mode
sweep, for the reason the next paragraph gives.

\paragraph{Caveat, stated because it is load-bearing.} The Fourier operator's mode
sweep is confounded with width at a matched budget --- the spectral weights cost
\(\mathcal{O}(w^2 m^2)\), so width falls as modes rise. The monotone degradation
across that sweep therefore \emph{cannot} be read as ``more Fourier modes do not
help''. The flat-in-\(p\) statement is made for DeepONet, where \(p\) and width are
controlled separately, and is not claimed for the Fourier operator.

\subsection{L6 --- separate identification of the kinetic and equilibrium objects}

This was the project's primary hypothesis, and its pre-registration was committed to
version control before the first arm trained, so its ordering is externally
verifiable from the repository history, as it is for the other rungs listed in
Section~\ref{sec:reproducibility}.

Both invalidation checks passed before any verdict was computed. The kinetic object
is unrecoverable from observable descriptors (\(R^2 = \nLsixKineticRtwo\), worse than
predicting the mean) while the equilibrium shape is recoverable
(\(R^2 = \nLsixStepRtwo\) for the step position, \nLsixQmaxRtwo\ for the capacity).
That asymmetry is exactly what the hypothesis needs.

\paragraph{Primary estimand: not supported, and a weak null.} The separate-versus-joint
difference shows no detectable dependence on Damk\"ohler number: slope \nLsixSlope\
per decade, interval [\nLsixSlopeLo, \nLsixSlopeHi], Pearson \(r = \nLsixSlopeR\),
over \nLsixDaDecades\ decades of the informative band, Damk\"ohler
\nDaMatMin--\nDaMatMax\ per material (Figure~\ref{fig:ladder}B). In the pre-declared
words, \textbf{the mechanism is not supported}. What the null cannot exclude must be
said with it: carried across the tested range, the upper end of the slope interval
allows a change of \nLsixSlopeSpanMax\ nRMSE, about \nLsixSlopeVsEffect\ times the
main effect below. A Damk\"ohler dependence as large as the benefit itself is
therefore not ruled out; what is ruled out is a dependence strong enough to have been
seen at this design's resolution. The axis is also not purely kinetic: for the
\nDcExtendedPct\,\% of runs whose horizon was extended, Damk\"ohler exceeds the kinetic
ratio \(k_{\mathrm{LDF}}t_{\mathrm{st}}\) by the extension factor
(Section~\ref{sec:methods}), so part of its spread is the length of the run.

\paragraph{The scope of that refutation, stated because it is the strongest attack
on it.} Kinetic control means Damk\"ohler well below unity, and this dataset never
goes there: its lowest per-material value is \nDaMatMin. The hypothesis predicted
that separating the kinetic object pays most where kinetics dominate, and we have
refuted it over a range whose slow-kinetics end is merely \emph{less fast}, not
slow. A design that sampled \(\mathrm{Da} < 1\) could in principle recover the
interaction. We did not sample it, and the reason is structural rather than an
oversight: the kinetic coefficient is not a free parameter in this dataset but is
derived from particle size through the Glueckauf relation, so
\(\mathrm{Da} \propto d_p^{-2}\) and reaching \(\mathrm{Da} = 1\) from our slowest
bed would take a pellet of about \nDaPelletForUnity\,\si{\milli\metre}. That is
\nDaPelletsAcross\ pellets across the \nDataBedLenCm\,\si{\centi\metre} column, where a
one-dimensional plug-flow model is no longer the right description --- the same
objection Section~\ref{sec:limitations} raises against our own anchor bed. Decoupling
\(k_{\mathrm{LDF}}\) from \(d_p\) to reach that regime would restore precisely the
physical inconsistency this dataset was built to remove. The refutation therefore
holds across the range a physically consistent bed of this kind can occupy, and is
not a statement about the kinetically-controlled limit.

\paragraph{Secondary estimand, significant, and a sign reversal.} Joint fitting
gives \nLsixJoint\ and separate identification \nLsixSeparate: separate is better by
\nLsixDiff, interval [\nLsixLo, \nLsixHi], consistent in all five folds and all
three seeds. On the earlier dataset separate was \emph{worse}, with an
interval spanning zero. The direction has flipped, and the earlier number is not
quoted here because it belongs to a design that could not test the estimand at all.

\paragraph{Why the test has power now, and its limit.} The between-material standard
deviation of the arm difference fell to \nLsixBetweenSD, or
\nLsixBetweenOverBase\,\% of the base error, from \nLsixLegacyBetweenOverBase\,\% previously --- a \(\nLsixNoiseDrop\times\) noise reduction from Sobol sampling and physically consistent
kinetics. With \nNclustFolds\ clusters instead of \nNclustLegacy, that is what turned an unanswerable
question into an answerable one. Power is \nLsixPowerFour\,\% at a 4\,\% effect and
\nLsixPowerTwo\,\% at 2\,\%, with size at the null \nLsixSizeNull\,\%; the minimum
detectable effect at 80\,\% power is \nLsixMDE\,\%. The observed effect sits just
under it, so the point estimate may be optimistic and \textbf{the lower bound is the
conservative reading}.

\paragraph{Verdict, in the pre-declared words.} Structural decomposition gives a
small but real and reproducible improvement in transfer to unseen materials, with no
detectable dependence on Damk\"ohler number. \textbf{The pattern is that of an
inductive bias rather than a kinetic identifier} --- which is not what the hypothesis
predicted --- though, as the slope interval shows, a kinetic component of the
benefit's own size is not excluded, so this reading is the better-supported one
rather than an established one. This is a more precise statement than either a null or a naive
positive, and it is available only because the mechanism was pre-registered as the
primary estimand and could therefore fail on its own terms.

\subsection{L7 --- is learning needed at all}

Scored like for like on the exit curve, because the closed forms predict a curve and
not a field. The learned arm's out-of-fold prediction reaches \nLsevenLearned{}
against \nLsevenKlinkenberg\ for the best closed form, \nLsevenShock\ for
equilibrium theory and \nLsevenPattern\ for constant-pattern theory. The best closed
form is \nLsevenRatio{}\(\times\) worse, with difference \nLsevenDiff\ and interval
[\nLsevenLo, \nLsevenHi]. \textbf{Verdict, in the pre-declared words: LEARNING
IS JUSTIFIED} on dataset v2 --- every classical form is significantly worse than the
learned reference on the exit curve.

The transferable statement is neither ``classical wins'' nor its reverse, but that
\textbf{a classical model wins exactly when the physics it assumes is the physics
that is there}. Where a problem is linear and time-invariant, a Green's function is the exact
solution operator and a linear convolution is the right model, not an approximation
to one. Adsorption with an inflected
isotherm is nonlinear, and the two-wave structure violates all three closed-form
assumptions: a linear isotherm, a non-spreading favourable front, and infinitely
fast kinetics.

%% ===========================================================================
\section{What binds the error}
\label{sec:mechanism}

\begin{figure}[t]
\centering
\includegraphics[width=\textwidth]{../figures/Fig5_mechanism.pdf}
\caption{The mechanism. \textbf{A}~Held-out error against basis size: the floor
falls \(\nFloorDrop\times\) while DeepONet and the pointwise regressors stay flat
(the Fourier operator has no basis size and is drawn as a level). \textbf{B}~The two terms of a linear
reconstruction, measured in the \emph{optimal} basis. \textbf{C}~Per-mode
\(R^2\) of the coefficient map on held-out materials.}
\label{fig:mechanism}
\end{figure}

\paragraph{In plain terms.} Most models here build a breakthrough field as a weighted
sum of stored template shapes (proper orthogonal decomposition plus a regressor, and
DeepONet); the Fourier operator learns the field directly. A template model can be
limited by its templates or by its recipe for weighting them from a material's
isotherm and pellet properties. It is the recipe: with perfect weights the templates
alone would reproduce the field far more accurately than any recipe we can learn,
and the most accurate learner by point estimate still sits \(\nFNOxFloor\times\)
above that level.

A linear reconstruction can fail in exactly two ways: the basis cannot represent the
field, or the map from parameters to coefficients cannot be learned. We separate
them in the \emph{optimal} basis, so that the answer cannot depend on the basis any
particular method happened to learn (Figure~\ref{fig:mechanism}).

The basis error falls from \nLfiveBasisSmall\ to \nLfiveBasisLarge\ as \(p\) goes
from 8 to 128. The coefficient-map error, with gradient boosting as the regressor
(the better of the regressors tried, at every basis size), moves from
\nLfiveCoefSmall\ to \nLfiveCoefLarge\ --- a change of \nLfiveCoefChangePct\,\%
across a sixteenfold change in basis size, against a basis error that falls by
\(\nFloorDrop\times\) over the same range. The
per-mode reason is that on held-out materials the map predicts few modes at all:
\nLfiveCountModes\ clear \(R^2 > 0.2\), the leading \nLfiveLeadModes\ and
\nLfiveIsolatedModes\ isolated later ones. The median \(R^2\) of the
\nLfiveTailModes\ modes after the leading run is \(\nLfiveTailMedian\), worse than
predicting the mean, and
\nLfiveTailFracNeg\,\% of them are negative.

Against an oracle reconstruction handed the true coefficients for its first \(k\)
modes, the best DeepONet found anywhere performs like a \nLfiveOracleONet-mode
oracle, and the best Fourier operator like a \nLfiveOracleFNO-mode one. Three
independent function classes --- proper orthogonal decomposition plus random forest,
plus gradient boosting, and DeepONet --- are all flat in rank while the floor falls by a factor of \nFloorDrop. Training basis and coefficients together, or learning the
whole field at once, is worth a constant factor --- up to \(\nOperatorLever\times\)
from the pointwise regressor to the Fourier operator on this encoding, a lever of
the same order as the learning curve's (the two ratios come from different designs,
and the operator ratios are point estimates between best-of-grid arms, with no
interval of their own)
--- and it does not close the distance to the
floor.

\paragraph{The decomposition holds on the full design.} The numbers above were
measured on the legacy held-out set, and the claim they carry is the one this paper
rests on, so we re-measured it where every other rung reports: \nLfiveVtwoNmat\
materials, five folds pooled out of fold, three seeds. The basis floor falls from
\nLfiveVtwoFloorSmall\ to \nLfiveVtwoFloorLarge, a factor of
\(\nLfiveVtwoFloorDrop\); the coefficient-map error moves from \nLfiveVtwoCoefSmall\
to \nLfiveVtwoCoefLarge, \nLfiveVtwoCoefChangePct\,\% of itself. That change is
statistically resolved --- interval [\nLfiveVtwoCoefLo, \nLfiveVtwoCoefHi] --- and
immaterial: the design is now powerful enough to see a difference far too small to
matter, and the ratio of the two errors grows with \(p\) exactly as before. At the
largest basis the map predicts about \nLfiveVtwoModes\ modes on held-out materials.

\paragraph{And this is scoped to a sample size.} An
earlier version of this section stated the mode count as a property of the map. It
is not: it is what a given number of training materials buys. Measured against
material count, modes clearing \(R^2 > 0.2\) grow from \nLCmodesSmall\ to
\nLCmodesLarge\ over 12 to 192 materials and are still climbing. Three counts of
modes clearing \(R^2 > 0.2\) appear in this paper, and they are the same statistic
under different conditions: \nLfiveCountModes\ on the legacy design, and
\nLCmodesLarge\ and \nLfiveVtwoModes\ on the larger one at the same number of training
materials, the first for the L1 perceptron on its own basis and the second for the
decomposition's regressor on the largest basis. They agree on what matters here ---
a learner predicts a modest set of mostly leading modes and little beyond it. The retraction is
A21, and it withdrew the replacement text for an earlier retraction --- the same
inference, made twice.

\paragraph{The \(n\)-width is real and was not what binds.} The residual
\emph{energy} of this snapshot family decays algebraically, as
\(k^{\nNwidthExpC}\) in the number of modes \(k\), with \(R^2 = \nNwidthRtwoC\) over
\(k = 8\)--64. The Kolmogorov \(n\)-width, whose \(n\) is this mode count, is the
residual \emph{norm}, the square root of that, so it decays as \(k^{\nNwidthExpErr}\)
and halving the error requires
\(\nNwidthHalving\times\) the modes permanently. We separate the two because an
earlier draft quoted the energy exponent and called it the \(n\)-width, next to a
halving factor computed correctly from the norm --- two numbers in one sentence
that disagreed by a factor of two in the exponent. Any method whose output is a linear
combination of a fixed basis is bounded below by it. Our frozen protocol called this
the load-bearing result of the ladder and predicted that DeepONet would track it.
The prediction was wrong: the bound is real, and it is inactive. Withdrawing it is
A17. For transport-dominated problems the slow decay is established --- Ohlberger
and Rave \cite{ohlberger2016reduced} prove a \emph{lower} bound of \(\tfrac12 k^{-1/2}\) for linear advection
with jump discontinuities, which is a bound and not a rate, and we state it that
way. Our measured decay is faster than that bound, and the reason is the
hypothesis: their fields carry jump discontinuities, ours are smoothed by axial
dispersion and by finite-rate kinetics, so the worst case their theorem
constructs does not occur here. That is also why the bound can be simultaneously
real and inactive --- the two statements this paragraph has to hold together.

%% ===========================================================================
\section{The coordinate change, and where the obstruction moved}
\label{sec:warp}

\begin{figure}[t]
\centering
\includegraphics[width=\textwidth]{../figures/Fig6_warp.pdf}
\caption{The two-wave frame. \textbf{A}~Modes for 99.9\,\% of training variance per
frame. \textbf{B}~The verdict: the honest arm ties the fixed frame, the
oracle arm is significantly better and is an upper bound.}
\label{fig:warp}
\end{figure}

\subsection{Scope and prior art, stated before the result}

Aligning multiple transports is \textbf{not new}, and none of it is claimed here.
The two-transport case is the founding worked example of the shifted proper
orthogonal decomposition \cite{reiss2018shifted}, where a pressure pulse splitting into two waves is
reconstructed with two modes --- one per co-moving frame --- against more than
eighty ordinary modes; the general formulation covers arbitrary numbers of
independent transports. Multi-front variants follow in several later works \cite{nair2019transported,mendible2020dimensionality,krah2025robust,krah2023front,zorawski2024neural}, and a neural-network shifted decomposition predates the last of these by two years \cite{papapicco2022nnspod}.
Shape--timescale decomposition of breakthrough curves specifically has been
published \cite{lee2026scaleup}, with a \emph{single} characteristic time and a
uniform rescaling.

Nor is the map new mathematics. Piecewise landmark registration --- a monotone warp
built by interpolating between identified landmark times --- is the standard
alignment method in functional data analysis \cite{kneip1992statistical}, with a
landmark-free alternative in Ramsay and Li \cite{ramsay1998curve}, and \textbf{our two-landmark linear
case reduces exactly to it}. Zucatti and Zahr \cite{zucatti2026model} use landmark-based registration for
reduced-order models of convection-dominated problems, which is the same primitive
applied to shock features in space rather than to fronts in time; it is the nearest
neighbour of this section and we name it as such. Registration-based reduction is
itself an established family \cite{taddei2020registration,taddei2026erratum}, so this
narrows our claim rather than restoring novelty.

\textbf{What is ours is the measurement.}

\subsection{One shift makes it worse; two make it much better}

The standard remedy for a transport-dominated \(n\)-width is to factor the front
position out before projecting \cite{reiss2018shifted}. We warped each field onto its own local arrival time
and swept the alignment level. \emph{It makes the \(n\)-width worse}: aligning on
the half-height crossing takes the modes needed for 99.9\,\% of training variance
from \nNwidthOriginal\ to \nNwidthSingle. And it does not help accuracy even given
the exact front trajectory for free --- handed the oracle warp, the co-moving frame
still loses to the fixed frame. The reason was measured before it was invoked: the separation between the two waves spans a factor of \nGapRange, and one shift cannot align
two waves whose separation varies that much.

Using \emph{two} shifts, one per wave, and normalising the interval between them (Figure~\ref{fig:warp}),
\begin{equation}
\sigma(z,t) = \frac{t - t_{\mathrm{lo}}(z)}{t_{\mathrm{hi}}(z) - t_{\mathrm{lo}}(z)},
\label{eq:warp}
\end{equation}
collapses the \(n\)-width where one shift inflated it: \nNwidthOriginal\ \(\to\)
\nNwidthTwoWaveA\ modes, or \nNwidthTwoWaveB\ for the wider landmark pair. A third
pair was \textbf{excluded} because a narrow window divides by a small span and its
interpolation floor reached \nComovingFloorPct\,\% of the fixed-frame error it is judged
against, where the guard's threshold is \(10\,\%\) --- the guard caught it, and it is
reported rather than dropped silently.

\subsection{The verdict, what it can resolve, and the relocation}

Three seeds across every stochastic component, both arms computed in the same run,
paired cluster bootstrap by material:

\begin{table}[t]
\centering
\caption{The two-wave frame. The oracle arm is handed the true front trajectories
and is \emph{not} comparable with any arm that is not; it is reported to size the
remaining headroom, never as a result.}
\label{tab:warp}
\resizebox{\textwidth}{!}{%
\begin{tabular}{lrl}
\toprule
arm & held-out nRMSE & verdict vs.\ the fixed frame \\
\midrule
fixed frame & \nWarpFixed & --- \\
two-wave, predicted fronts & \nWarpPredicted & \nWarpPredDiff, CI [\nWarpPredLo, \nWarpPredHi]: no difference \\
two-wave, oracle fronts & \nWarpOracle & \nWarpOracleDiff, CI [\nWarpOracleLo, \nWarpOracleHi]: significant \\
\bottomrule
\end{tabular}}
\end{table}

Two statements, and both are needed. The honest arm does \emph{not} beat the fixed
frame. But the oracle gap is significant and large: given the true front
trajectories, the same proper-orthogonal-decomposition-plus-regressor arm is
\nWarpHeadroomFixed{}\(\times\) better \emph{than the fixed frame}, and
\nWarpHeadroom{}\(\times\) better than the same arm with predicted fronts. The two
denominators differ and we name which is which, because the quantity a reader
wants --- what the coordinate change would buy if the fronts were located
perfectly --- is the first.

\textbf{So the gain is not recovered once the fronts are predicted}, and this is
a measured null, not a demonstrated absence. On \nWarpNmat\ held-out materials the
interval for the predicted-front arm admits anything from a \nWarpPredImproveMax\,\%
improvement to a \nWarpPredDegradeMax\,\% degradation of the fixed frame; it rules
out the oracle's gain, not a modest one. With that scope, this is a relocation of the
obstruction rather than only a null. Through L1--L5 the obstruction was the
parameter-to-coefficient map, at the number of modes a given training set supports. It
is now ``two smooth monotone one-dimensional curves are not located accurately
enough'' --- a better-posed problem, with a target a front-locating model can be held
to. The oracle's gain is not unique to this route: the Fourier operator of L5 reaches a
comparable error on the same held-out materials without any oracle.

\paragraph{Monotonicity is not the lever.} The obvious first attack is that the
front predictor does not know the constraints the true trajectories satisfy exactly.
It does violate them, in \nWarpNonmonoLo\,\% of samples at the lower landmark, while the true trajectories
violate them in none. But projecting onto the constraint set changes \(R^2\) by
at most \nWarpMonoMaxDRtwo{} and makes the reconstruction significantly
\emph{worse}: \nMonoRawNrmse\ against \nMonoProjNrmse, a difference of
\nMonoDiff\ with interval [\nMonoLo, \nMonoHi]. That interval is much tighter than
Table~\ref{tab:warp}'s, and the reason is pairing rather than power: the projected
arm is a deterministic transformation of the raw arm's own predictions, so almost
all of the between-material variance cancels, whereas the two-wave and fixed frames
are separately fitted arms. A smaller difference can therefore be resolved here
than there, and the two intervals are not comparable. The violations are frequent
but negligible in magnitude. A monotone-by-construction front predictor is not worth
building; the remaining headroom is raw predictive accuracy of the front predictor,
which more training materials or a better learner could supply.

\paragraph{A constant-pattern objection, pre-empted.} A referee may observe that
constant-pattern behaviour --- a front propagating without changing shape --- is
long established in fixed-bed adsorption \cite{knox2016limitations}, and that equation~\eqref{eq:warp} merely
re-derives constant-pattern scaling. It does not, and the reason is the same
measurement that killed the single-front frame: when the two fronts move at
different rates, and their separation varies \nGapRange-fold, no single scale factor
aligns both. That is precisely the regime a one-parameter dimensionless time cannot
handle, and it is why the second landmark is doing the work.

%% ===========================================================================
\section{What the surrogate costs, and when it repays}
\label{sec:cost}

A surrogate that is never asked enough questions is more expensive than the solver it
replaces, and the crossing point is a number, not an attitude. The introduction argued
that a synthetic study can only measure accuracy per unit cost; this section pays that
debt, because a ladder of accuracy verdicts with no cost axis is only half an argument
for a screening tool.

\paragraph{The unit.} All costs are given in \emph{solve-equivalents}: multiples of one
reference simulation at production settings. That removes the machine, the thread count
and the wall clock from the comparison and leaves a ratio a reader can carry to their own
hardware. On ours, the \nCostNSolves\ accepted simulations of the dataset took
\nCostGenWall\ hours on \nCostWorkers\ workers, so one solve-equivalent is
\nCostSolveSec\ worker-seconds and the whole training set cost about
\nCostGenCoreHours\ core-hours. We report worker-seconds rather than
core-seconds and say why: a worker's own numerical libraries may use more than one
thread, so this is an upper bound on the serial cost of a solve and a lower bound on
what it occupies of a machine.

\paragraph{The accounting.} Against \(N_q\) screening queries the reference costs \(N_q\)
solve-equivalents. The surrogate costs its training set, plus its fit, plus \(N_q\) times
whatever a query costs, so the two break even at
\begin{equation}
N_q^{\mathrm{be}} \;=\; \frac{N_{\mathrm{sim}} + c_{\mathrm{fit}}}{1 - r},
\qquad r \;=\; \frac{\text{cost of one query}}{\text{cost of one solve}} .
\label{eq:breakeven}
\end{equation}
with \(N_{\mathrm{sim}}\) the number of training simulations and \(c_{\mathrm{fit}}\)
the fit, both in solve-equivalents. Fitting the reported arm takes \nCostTrainSec\ seconds, median over
\nCostTrainSeeds\ recorded fits --- \nCostTrainEq\ solve-equivalents, which is noise
next to the data. And because \(r \ll 1\) --- a forward pass through three
256-unit layers and three 64-mode reconstructions, against an implicit solve on
\nGridNz\ cells carried to breakthrough --- equation~\eqref{eq:breakeven} is bounded
below by \(N_{\mathrm{sim}} + c_{\mathrm{fit}}\) \emph{without measuring \(r\) at
all}. We report the bound, because a bound that needs no timing cannot be argued with
by pointing at the hardware. Measured as well, on an idle machine, a full-field
query takes \nCostQueryMs\ ms, so \(r = \nCostQueryRMant \times 10^{\nCostQueryRExp}\) and a query is about
\nCostQuerySpeedup\ times cheaper than a solve, which makes the bound and the exact
break-even indistinguishable. That ratio is per query at the surrogate's stated
held-out error, not a claim that the surrogate replaces the solver.

\begin{table}[t]
\centering
\caption{The accuracy--cost frontier of the reported arm, in solve-equivalents. Each
row is a point on the learning curve of Section~\ref{sec:ladder}, priced. The last
column is a \emph{lower} bound on the number of screening queries at which the
surrogate becomes cheaper than the reference, and it holds whatever a query costs.}
\label{tab:cost}
\resizebox{\textwidth}{!}{%
\begin{tabular}{rrrr}
\toprule
training materials & training simulations & held-out nRMSE & break-even queries \(\geq\) \\
\midrule
12  & \nCostSmallSims & \nLCsmallest & \nCostSmallBreakEven \\
24  & \nCostSimsB     & \nCostErrB   & \nCostBeB \\
48  & \nCostSimsC     & \nCostErrC   & \nCostBeC \\
96  & \nCostSimsD     & \nCostErrD   & \nCostBeD \\
192 & \nCostLargeSims & \nLClargest  & \nCostBreakEven \\
\bottomrule
\end{tabular}}
\end{table}

\paragraph{What Table~\ref{tab:cost} says.} Three things, and the third is the one that matters.

First, the surrogate at its best accuracy cannot repay its own construction in fewer
than \nCostBreakEven\ screening queries. Below that the solver is both cheaper
\emph{and} exact, and building a surrogate is a mistake. Above it --- and a materials
screen over a framework database is far above it --- the surrogate is the cheaper
option, at the accuracy in the table's last row. The solver still finishes: at
\nCostSolveSec\ worker-seconds a solve, a large screen is a matter of core-hours, not
of feasibility, and whether a coarsened solver at matched accuracy narrows that gap
is a baseline not yet run. This is not a criticism of the approach; it is the condition under which
the approach is the right one, stated in the units the decision is actually made in.

Second, the learning curve is a price list. Going from 12 training materials to 192
multiplies the data cost by \nCostDataRatio{}\(\times\) and divides the error by
\nLCratio{}\(\times\) --- which is the exponent \(N_{\mathrm{mat}}^{-\nLCbeta}\) of
Section~\ref{sec:ladder} restated as money. A practitioner who needs
\nLCsmallest\ nRMSE should buy \nCostSmallSims\ simulations and stop. The list is
priced for one fixed learner, the L1 perceptron; L2 found a configuration
\(\nLtwoVsLonePct\,\%\) better at the largest size, and the operator comparison of
Section~\ref{sec:mechanism} a larger lever still, so a better learner shifts every row
down and the frontier with it. The exchange rate between data and error is the
transferable part, not the absolute rows.

Third, and this is why the negative results in this paper are worth their space: every
rung that was \emph{eliminated} is an intervention that would have added cost without
buying accuracy. A richer per-edge basis (Kolmogorov--Arnold against perceptron) and
the basis size of a linear-reconstruction operator cost fit time and nothing else, and the closed forms are cheaper but worse.
Not everything was eliminated, and the levers that remain are not equal. Training
materials are the one with a price tag attached, and the frontier is where that price
is paid. Capacity is not eliminated, but a better configuration found by search costs
fit time, not data. The residual as a loss is not priced here: its rung is being
re-run with the corrected equation (B72). The choice of
operator moves the accuracy without moving the data cost: on the legacy held-out
materials the Fourier operator reaches \nLfiveFNO, against \nWarpOracle\ for the
coordinate change of Section~\ref{sec:warp} even when that change is handed the true
fronts, so the warp is no longer the only route on the table that improves accuracy
at a fixed data cost. This is the one place we compare absolute errors across rungs,
and it is licensed: L5 and the coordinate change share the legacy split, the same
subsampled concentration field and the same per-sample normalised error (each
sample's root-mean-square error over its own range); the subsampling grid is built
by the same rule in both codes, which implement it separately.

\paragraph{What this is not.} It is not a claim that the surrogate and the reference are
interchangeable. They are not: the reference answers exactly and the surrogate answers
with the held-out error in the third column. Every statement here is cost \emph{at a
stated accuracy}, and the accuracy travels with the cost in every row.

%% ===========================================================================
\section{Corrections and retractions}
\label{sec:corrections}

This section is not an appendix. \nLedgerA\ claims made in this project have been
withdrawn, and \nLedgerB\ further defects were caught and corrected --- most, but
not all, before they reached a reported number. B43 is the exception and is named
as one: three quoted comparisons, including the quantitative basis of a retraction,
had no script behind them at all, and the numbers had already been written down
before anyone could check them. \nLedgerOwn\ of the \nLedgerTotal\ entries are marked ``our own error'' and
several more are self-attributed in other words --- in the analysis, in the
validation gates, in the power simulations, in the citation ledger itself, and in
the frozen protocol.

We report this for two reasons. The first is that the corrections are load-bearing:
four of them reversed headline verdicts, and a reader who does not know which
numbers moved cannot evaluate the ones that did not. The second is evidential.
Disclosed self-correction is received well, and the Loss-of-Confidence Project
documents that \cite{rohrer2021loss}; the alternative --- a paper in which nothing ever went wrong --- is
less credible, not more.

\paragraph{The classes that recur.} Five patterns account for most of the ledger.
\begin{enumerate}
\item \textbf{A property measured at one sample size, stated as intrinsic.} A18,
A21 and A26 are the same error three times: the learning-curve verdict, the
coefficient-map mode count, and the capacity optimum. Each was true at the size it
was measured at and false as a general statement.
\item \textbf{Normalising by a quantity that can vanish.} A15 destroyed every
time-axis number in one rung when a window-local range collapsed after breakthrough;
B34 and B41 repeated it in diagnostics; and the script written for this paper to
check the isotherm did it again in its first version, dividing by a ceiling that
underflows to zero.
\item \textbf{Comparing against a weak configuration of the competitor.} B14, B22,
B23 and A20 --- an unscaled target vector, a single learning rate that suits one
family, a two-seed selection rule, and a recorded instruction to update a table that
was never carried out.
\item \textbf{Undertrained arms scored as architectures.} A19 and B24: \nBtwentyfourUntrained{} of
\nBtwentyfourRuns{} runs in one rung reported the error of a random initialisation, because
an aggressive patience does not handicap architectures equally.
\item \textbf{Power and calibration.} A22 withdrew a claim that a null was well
powered when it had 7\,\% power; A25 corrected the correction in the opposite
direction; A13 withdrew a positive verdict that did not survive the calibrated
level.
\end{enumerate}

\paragraph{Four that reversed a headline.} The physics-as-a-loss rung reversed
twice. Its legacy claim that the residual \emph{hurts} did not survive a sweep of
weightings and the edge rule (B69); then every one of its verdicts was withdrawn as a
statement about physics when the residual itself proved to be the wrong equation (B72,
A28) --- the largest-scope withdrawal in the ledger, and the only one that voids a
whole rung. A17 withdrew the explanatory role of the
Kolmogorov \(n\)-width, which the frozen protocol had called the load-bearing result
of the ladder --- the bound is real and inactive. A22 turned a null we had described
as decisive into one that excludes only very large effects. A25 then corrected A22's
replacement in our favour, and is recorded on the same terms; it qualifies A22 rather
than reversing a fifth headline. The last internal review added B75, which removed
the word ``significantly'' from the Fourier operator's advantage over DeepONet once
the choice of both arms on the test materials was paid for. We do not count it among
the four: the direction of that comparison stands, and only its significance went.

\paragraph{What remains open.} One entry is deliberately unresolved. The classical
control in rung L7 uses an approximation attributed in our bibliography to
Klinkenberg's 1948 paper \cite{klinkenberg1948numerical}, while the standard
secondary source attributes that equation to a different Klinkenberg paper of 1954
\cite{klinkenberg1954heat}. Neither is retrievable
electronically, and the ambiguity is recorded rather than settled by choosing the
year that looks right.

%% ===========================================================================
\section{Limitations}
\label{sec:limitations}

\paragraph{Best-of-grid selection on the test materials.} Several comparisons pit the
best of a grid --- of learning rates, capacities, weightings or operator settings ---
against a single reference, with the best chosen on the same held-out materials the
comparison is scored on. A per-pair interval is valid for a pair fixed in advance and
anti-conservative for one chosen after looking. For these we report a studentised
max-statistic cluster bootstrap, the single-step form of the procedure of Romano and
Wolf \cite{romano2005stepwise}: materials are resampled jointly across all candidate
arms, each arm's deviation is scaled by its own standard error, and the critical value
is the quantile of the largest scaled deviation. The resulting interval holds for
every candidate at once, so for the selected one however it was selected; on
simulated families with no true effect it holds its level where the naive test on the
winner does not. The worked case is L2: its best configuration, chosen from
\nLtwoSelK\ on the held-out materials, keeps a gain over the L1 setting of
\(\nLtwoSelLo\)--\(\nLtwoSelHi\,\%\) once the choice is paid for. The same procedure
will be applied unchanged to the corrected L4 re-run (on the withdrawn residual the
critical value was \nLfourbSelQ\ standard errors, wider than a single pair needs).
The correction changes one verdict. The Fourier operator's advantage over DeepONet
rests on two choices made on the test materials, one per family. Paid for over all
\nLfiveSelK\ pairs, it is no longer significant (Section~\ref{sec:ladder}, B75).
In L3 both families are taken at their best learning rate. Paid for in the same way,
\nLthreeSigSim\ of the \nLthreePairs\ comparisons remain significant, and none reverses
(Section~\ref{sec:ladder}). The cleaner remedy for every such choice --- selecting on a validation split
inside the training materials --- is the protocol change we recommend for any rerun.

\paragraph{One partial differential equation family.} Everything here is measured on
one model of one process. The ladder is a method; its verdicts are not claimed
beyond adsorption columns with inflected isotherms.

\paragraph{Solver-versus-experiment and surrogate-versus-solver are different error
budgets, and we keep them apart in those words.} Section~\ref{sec:anchor} is the only
solver-versus-experiment comparison in this paper: one condition, on a bed whose
aspect ratio is \nAnchorAspect, where a one-dimensional plug-flow model is at the edge of its
validity. Every other number is surrogate-versus-solver.

\paragraph{The dispersion closure needs a more careful statement than we first gave
it, and we have measured how much.} Our datasets use
\(D_L = 0.7 D_m + 0.5 d_p u\), with \(u\) the interstitial velocity. The coefficients are attributable to Wakao and
Funazkri \cite{wakao1978effect} rather than to the textbook we originally cited, and
Perry's handbook \cite{levan2019adsorption} presents that form as a \emph{lower bound} for \emph{nonporous}
particles, noting that the leading coefficient rises to as much as
\(20/\varepsilon_t\) depending on the intraparticle mass-transfer mechanism and that
for strongly adsorbed species the effective axial dispersion is much larger. Our beds
are porous, strongly adsorbing frameworks taking up water --- the case that statement
excludes.

That form is also the high-particle-P\'eclet \emph{limit} of Edwards and Richardson's
correlation \cite{edwards1968gas}, whose mechanical coefficient carries an extra
factor \(\left[1 + 9.7\,D_m/(d_p u)\right]^{-1}\). Whether the limit applies is a
question about our own parameter space, so we measured it rather than assuming it.
On the interstitial velocity the particle P\'eclet number \(d_p u / D_m\) runs
\nPePartMin--\nPePartMax\ across the dataset, median \nPePartMed, spanning
\nPePartDecades{} decades. Over that range the two published forms disagree by a
median factor of \nDispRatioMed, exceeding \(1.1\times\) for \nDispFracAboveTenPct\,\%
of samples and \(1.5\times\) for \nDispFracAboveHalf\,\%, with a worst case of
\nDispRatioMax. The disagreement is bounded and modest for the dataset --- but it is
not modest for the anchor.

\begin{figure}[t]
\centering
\includegraphics[width=\textwidth]{../figures/Fig7b_anchor_posthoc.pdf}
\caption{Post-hoc dispersion sensitivity on the anchor bed, labelled as such and
reported as a limitation rather than as a result. The pre-registered primary uses the
high-P\'eclet limit of the correlation; the second curve is the \emph{same}
correlation without that truncation, at the same particle size, velocity and
molecular diffusivity; the third is the authors' own Bruggeman closure. The measured
95\,\% arrival time is bracketed by the last two.}
\label{fig:anchorposthoc}
\end{figure}

The anchor bed sits at a particle P\'eclet number of \nAnchorPePart, and the
disagreement between the two forms is maximised at \nDispWorstPe, where it reaches
\nDispWorstRatio\(\times\). In other words the one bed in this paper that is compared
against an experiment is also the one place where the truncation matters most: there
our closure disperses \nAnchorDispRatio\(\times\) more than the correlation it is the
limit of, taking the column P\'eclet number from \nAnchorPeColER\ to \nAnchorPeCol.
Re-solving that bed under the untruncated form --- a labelled post-hoc sensitivity,
not a fit, and not the pre-registered result --- moves the 95\,\% arrival time from
\nAnchorTninetyfiveModel\ to \nAnchorTninetyfiveER\ minutes against
\nAnchorTninetyfiveExp\ measured, the plateau from \nAnchorPlateauModel\ to
\nAnchorPlateauER\ against \nAnchorPlateauExp, and the error against the digitised
effluent from \nAnchorNrmse\ to \nAnchorNrmseER\ (Figure~\ref{fig:anchorposthoc}).
The authors' own Bruggeman closure goes further, to \nAnchorNrmseBrug, but overshoots
in the other direction at \nAnchorTninetyfiveBrug\ minutes; the measured value is
bracketed by the two closures. \textbf{So a substantial part of the anchor's second
named discrepancy is not the solver and not the bed geometry, but one term evaluated
at its limit}: removing the truncation removes \nAnchorTcutPct\,\% of the model's
excess 95\,\% arrival time over the measurement, and \nAnchorNrmseCutPct\,\% of its
error against the digitised effluent. Those two shares differ because the error is
dominated by the first-wave timing, which no dispersion closure touches --- which is
itself the reason to quote both rather than an average of them. None of this affects
a surrogate result, because the solver is the reference
whatever closure it uses and every arm sees the same reference; it bounds how far the
reference may be read as physical, which is exactly what a limitation should do.

\paragraph{The isotherm's low-humidity branch is wrong, and we know by how much.}
The Henry fraction was fitted to the step position and the capacity, never to
dynamics. The anchor comparison is the first measurement that constrains it and
shows the model holds too little water below the step. Separately, the step position
in real MOF-303 shifts with temperature \cite{hastings2024high},
while ours is fixed.

\paragraph{Part of the heat range is occupied by no placed framework.} The sampled
heat of adsorption runs from \nDcDHLo\ to \nDcDHHi\,kJ\,mol\(^{-1}\); the
water-harvesting frameworks we placed (Figure~\ref{fig:watermofs}) report \nWmofHeatMost\ to
\nWmofHeatLeast. Near its least exothermic edge the sampled range gives an isosteric
heat (\(-\Delta H + RT\)) below water's latent heat of vaporisation at room
temperature, so that water is bound more weakly than in its own liquid. That is the
regime of hydrophobic adsorbents rather than of the frameworks this paper is
motivated by. The range was pre-registered and is kept; verdicts averaged over it
weight that region as heavily as the occupied one.

\paragraph{Not every sample reaches the step.} A harvester is operated above its
step by design; the sampled space is not. \nChemBelowStepPct\,\% of the accepted
samples of the larger dataset (\nChemBelowStepLegPct\,\% of the legacy one) are fed
below their material's step humidity, so the feed never reaches pore filling and the
breakthrough is the primary-site front alone. These samples are in every verdict; a
surrogate built only for harvesting conditions would not need them.

\paragraph{The step does not move between feed temperatures.} The cluster affinity is
set per run so that the step sits at the material's step humidity at the feed
temperature (Section~\ref{sec:methods}). Within a run the step still moves with the
local bed temperature through the heat of adsorption, but across runs a material's
step humidity at the feed temperature is the same by construction. In a real
framework it may shift with temperature, by an amount we have not taken from the
literature. The datasets therefore do not test whether a surrogate transfers
the temperature dependence of the step.

\paragraph{Two parameters are not measurements.} Pellet density and bed porosity
are packing values, not MOF-303 measurements \cite{hanikel2021evolution}. No citable measured heat capacity for
MOF-303 was found, so it is reported as a sensitivity band rather than as a cited
constant.

\paragraph{No covalent-organic-framework case exists.} We searched and did not find
a published COF breakthrough measurement suitable for this comparison. The gap is
closed by saying so, not by inventing one, and the title claims only frameworks of
the class we actually model.

\paragraph{Three rungs are measured on the smaller dataset.} L3, L5 and the
coordinate change are measured over \nNclustLegacy\ held-out materials, while L1, L2, L6 and
L7 are measured over 240. Given that A26 exists precisely because a verdict measured
at 48 training materials did not survive 192, those three verdicts must
be read as scoped to the material count at which they were taken. The ordering of L3
has the same sign at every budget; it is large against the Chebyshev family, and
against the Gaussian radial-basis family at the two larger budgets it is not
significant once selection is paid for. The oracle gap of the
coordinate change is large; the other headline statements of these rungs are nulls,
whose resolution is the weakest in the paper at this design. Re-measuring them on the larger dataset is the first item of further work, and
for L5 it is queued under its own pre-registration. For the coordinate change it is
not possible as designed, and the reason is itself a result: the pre-registered
landmark screen fails on every fold of the larger dataset, because its fronts do not
have the shape the two-wave frame aligns. The Henry wave is not resolved at the
stored time resolution, and the second front is a gradual, dispersive rise rather
than a shock, so even a landmark placed relative to each cell's own
plateau falls within two samples of the start in at least \nWarpScreenBoundaryMin\,\%
of cells on every one of the \nWarpScreenNfolds\ folds. The coordinate-change result is
therefore scoped to the smaller design, where the fronts are resolved, and we do not
redefine the landmark until a screen passes.

%% ===========================================================================
\section{Conclusions}
\label{sec:conclusions}

We asked what binds the error of a surrogate for adsorption column dynamics when it
must predict a material it has never seen, and we asked it as a ladder of
pre-registered falsification tests rather than as a search for the best-fitting
network. The answer is assembled here in one place, because each piece of it was
earned on a different rung.

\paragraph{What does not bind.} Capacity is not eliminated under the pre-registered
rule, because a random forest is still improving at its capacity ceiling; the
\nLtwoNUshaped{} neural and boosting families saturate. A richer per-edge basis is not
the lever either --- at matched parameters a perceptron beats both
Kolmogorov--Arnold families on held-out materials in every comparison by point
estimate, \nLthreeSigSim\ of \nLthreePairs\ significantly after selection and none
reversed, so the premise that motivated this study is refuted on our own data. Basis size is not the
lever: the proper-orthogonal-decomposition floor falls \(\nFloorDrop\times\) while
DeepONet stays flat, and a Fourier neural operator, the remedy the approximation
theory names, is more accurate --- not significantly, once its selection is paid for
--- and still sits \(\nFNOxFloor\times\) above that floor. The
Kolmogorov \(n\)-width is real, and it is inactive; our own frozen protocol named it
as the load-bearing explanation and was wrong (A17). Structural decomposition of the
equilibrium and kinetic objects helps, by \nLsixDiff\ nRMSE, and the benefit shows no
detectable dependence on Damk\"ohler number, which reads as an inductive bias rather
than a kinetic identifier --- a weak null, whose interval admits a dependence of the
benefit's own size.
The physics-residual rung returns no verdict on physics: its residual advected the
gas at the wrong velocity (B72), so its results measure a mis-specified prior and are
withdrawn as statements about this column (A28) until the corrected re-run reports.

\paragraph{What binds.} The map from material parameters to the field, not the
basis it is expressed in. Error still falls as \(N_{\mathrm{mat}}^{-\nLCbeta}\) with the number of
training materials at the largest size we can afford, and for the learner of L1 the
whole of that gain is in the map: a basis refitted on more materials buys nothing
once the regressor sees the same data. How well the map is learned also depends on
the learner, and that lever is not small. On one encoding, split and metric, the best
pointwise regressor on proper-orthogonal-decomposition coefficients reaches
\nLfiveRankPodReg, a Kolmogorov--Arnold operator \nLfiveRankOKAN, DeepONet
\nLfiveRankONet\ and a Fourier neural operator \nLfiveRankFNO\ --- a factor of
\(\nOperatorLever\) from first to last, of the same order as the
\(\nLCratio\times\) that sixteen times more materials buys a fixed learner (the first
measured on the legacy design, the second on the larger one). What no
learner escapes is the distance to the floor: the most accurate by point estimate still sits
\(\nFNOxFloor\times\) above the linear-reconstruction floor at \(p = 128\), a reference
level it is not bound by. A change of coordinates
that aligns the two fronts collapses the \(n\)-width and would buy
\nWarpHeadroomFixed{}\(\times\) if the fronts were known, and nothing detectable when
they are predicted --- it relocates the obstruction, from a two-dimensional field to
two one-dimensional curves that must still be predicted from the same parameters.

\paragraph{What this means for a practitioner.} For a screening surrogate over
unseen adsorbent parameter sets, two kinds of effort pay: more, and more diverse, training
materials, and a learner that maps parameters to the whole field rather than to
coefficients one by one. A larger basis and a different nonlinearity at matched size
do not; whether a correctly specified physics residual does is the open question the
re-run of L4 answers. The cost accounting of
Section~\ref{sec:cost} gives the exchange rate for the first.

\paragraph{What is scoped.} Three verdicts --- L3, L5 and the coordinate change ---
were measured on \nNclustLegacy\ held-out materials rather than \nNclustFolds, and must be read at that
material count until they are re-measured. The reference solver reproduces the
published MOF-303 bed's two-wave shape and its half-breakthrough time with nothing
fitted, and misses its first-wave timing; of the excess in its late arrival time,
\nAnchorTcutPct\,\% is one dispersion term evaluated at a limit this bed does not
satisfy.

\paragraph{What we got wrong.} \nLedgerA\ claims were withdrawn and \nLedgerB\
defects caught in the course of this work, and they are reported in
Section~\ref{sec:corrections} on the same terms as the results. Several of the
paper's headline verdicts exist in their present form only because an earlier form
of them was falsified by a larger measurement. We regard that ledger as a result.
Every code cited in this paper (A1--A\nLedgerA, B1--B\nLedgerB) is given in full in
Supplementary Information S1, which is generated from the ledger file itself at every
build so that the two cannot disagree.

%% ===========================================================================
\section{Reproducibility}
\label{sec:reproducibility}

\paragraph{Which pre-registrations are verifiable, and which are self-attested.}
The pre-registrations for rung L6, for rungs L1/L2/L7, for rung L4b and for the
experimental anchor were committed to version control \emph{before} the first arm of
each was trained, and the ordering is externally checkable from the repository
history. The original protocol covering L0--L7 was not; it is self-attested, and we
say so rather than presenting all of them as equivalent.

\paragraph{Gates.} A verification harness of \nGates\ gates runs before any number is
reported, and no number from a failing category may enter this manuscript. The gates
encode defects that actually occurred: a plotting script that labels a model curve
must load a checkpoint or read a recorded results file and demonstrably compute
nothing itself; a training run whose best step is zero is a hard failure; a physics
reconstruction must be handed the dataset's own parameter keys with no default.

\paragraph{No result in this manuscript is typed by hand.} Every measured quantity
above is a macro resolved at build time from a named results file through a declared
path, or a ratio of such macros with its formula declared. A key whose file is
missing, whose path does not exist, or whose value is not finite is a build failure.
The remaining literals are design constants (grid values, fold and seed counts),
thresholds that define a quantity, citation years, and figures quoted from the
ledger's own history; each has a written reason and a recorded number of occurrences,
and a new occurrence refuses the build until it is justified again. Numbers written
as words are held to the same rule. The build checks text, so it cannot prove that a
literal is not a result; it can only make every one of them visible and argued for. This exists because defect B43 found three quoted comparisons
--- one of them the basis of a retraction --- that had no script behind them at all.

\paragraph{Consistency across runners, asserted rather than assumed.} The
configuration used as L1's arm appears twice more in L2's capacity sweep, trained by
a different runner on the same folds, basis and seeds. Their per-sample held-out
scores agree to \nLtwoIdentity. That is the strongest available evidence that
the rungs reported over ``the same folds'' really are.

\paragraph{Citations.} Nothing is cited until it has been retrieved, not searched.
That gate has caught a fabricated co-author, a swapped arXiv identifier between two
papers named in the same sentence, a publisher-supplied identifier pointing at a
different paper, a missing co-author, two references that do not contain the claim
they were cited for, and a correlation attributed to the wrong source. References
that could not be retrieved at the primary, and what each was read through instead, are
listed below; the list is generated by the same script that decides which references
are emitted, so it cannot disagree with the bibliography.
% GENERATED BY paper/references.py from VIA -- DO NOT EDIT.
\begin{itemize}
\item \cite{danckwerts1953continuous}: the boundary conditions, read through Pearson~\cite{pearson1959note}.
\item \cite{glueckauf1955theory}: the linear-driving-force factor, read through Perry's handbook~\cite{levan2019adsorption} and Cruz, Magalh\~aes and Mendes, who attribute the plain factor to the companion paper~\cite{glueckauf1947theory}, which is why both are cited.
\item \cite{klinkenberg1948numerical}: the breakthrough approximation, read through Seader, Henley and Roper's \emph{Separation Process Principles}, which attributes it to the later paper~\cite{klinkenberg1954heat}; neither primary could be retrieved, so both are cited.
\item \cite{anzelius1926uber}: the Anzelius--Schumann solution, read through Perry's handbook~\cite{levan2019adsorption}.
\item \cite{schumann1929heat}: the Anzelius--Schumann solution, read through Perry's handbook~\cite{levan2019adsorption}.
\item \cite{dodo2000model}: the two-term functional form, read through Buttersack~\cite{buttersack2019modeling}, who proposes a competing isotherm.
\end{itemize}


%% ===========================================================================
\section*{Data and code availability}
The solver, every runner and analyser, the verification harness, the
pre-registrations, the correction ledger and the citation ledger are in the project
repository \textbf{[public URL and archived DOI to be added at submission]}. Dataset
manifests are tracked; the field arrays are regenerable from the generator with the
recorded seed.

\section*{Declaration of competing interest}
The author declares no competing interests.

\section*{Funding}
This research received no specific grant from
any funding agency in the public, commercial or not-for-profit sectors.

\section*{CRediT authorship contribution statement}
\textbf{Abdulhamid Batayhi:} Conceptualization, Methodology, Software, Validation,
Formal analysis, Investigation, Data curation, Writing -- original draft, Writing --
review \& editing, Visualization, Project administration.

\section*{Declaration of generative AI and AI-assisted technologies}
During this work the author used Claude (Anthropic), an AI assistant, extensively:
to write and test code for the solver, the runners, the analysers and the
verification harness; to draft and revise the pre-registrations and the manuscript
text; to retrieve and check bibliographic records; and to act as an internal
reviewer. Every number in the manuscript is produced by a script from a results file
(Section~\ref{sec:reproducibility}), and the correction ledger records defects
whether they originated with the author or with the assistant. The author reviewed
and edited the content and takes full responsibility for the publication.

\bibliographystyle{elsarticle-num}
\bibliography{references}

\end{document}
